%%%%%%%%%%%%%%%%%%%%%%%%%%%%%%%%%%%%%%%%%%%%%%%%%%%%%%%%%%%
%%%%%%%%%%%%%%%%%%%%%%%%%%%%%%%%%%%%%%%%%%%%%%%%%%%%%%%%%%%
%%%%%%%%%%%%%%%%%%%%%%%%%%%%%%%%%%%%%%%%%%%%%%%%%%%%%%%%%%%
%%%%%%%%%%%%%%%%%%%%%%%%%%%%%%%%%%%%%%%%%%%%%%%%%%%%%%%%%%%
%%%%%%%%%%%%%%%%%%%%%%%%%%%%%%%%%%%%%%%%%%%%%%%%%%%%%%%%%%%
\chapter{[VIT] Vitalfunktionen}
\label{chp:VIT}
\AasRsN{RS~V}


\emph{Maintainer: Sebastian Gabriel}

% \AuthNotes{
% Changelog:
% 
% \begin{description}
%     \item[2009-02-25] Start Changelog. Kapitel mit Korr
%     Koch korr. 
%     \item[2009-mm-dd] in \LaTeX übernommen
%     \item[2009-04-19] anhand der OpenOffice-Version überarbeitet.
% Bilder fehlen noch. 
% \item[2009-05-06] GAB: Abschnitte "`Blutdruck"' und
% "`Atemvolumina"' aus Anatomie-Teil übernommen.
% \end{description}
% 
% }

\minitoc

\clearpage

\paragraph{Überblick}~

\begin{table}[H]
\Captionof{table}{Vitalfunktionen.}

\begin{tabularx}{\linewidth}[H]{XX}
\TabHeading{Vitalfunktionen I. Ordnung}
 &
\TabHeading{Vitalfunktionen II. Ordnung}
\\\addlinespace\midrule\addlinespace
\begin{itemize}
  \item Bewusstsein
  \item Atmung
  \item Kreislauf
\end{itemize}
\Fazit{Wichtigste, grundlegende Funktionen des Körpers

Ausfall einer dieser Funktionen bedeutet immer akute
Lebensgefahr und muss dringlich behandelt werden!}
 &
 \begin{itemize}
   \item Wasser-/Elektrolythaushalt
   \item Säure-/Basenhaushalt
   \item Temperaturhaushalt
   \item Hormonsystem
   \item Stoffwechsel
   \item Immunsystem
 \end{itemize}
\\\addlinespace\bottomrule
\end{tabularx}
\end{table}


\Layout{2}{NACA -- Score}
{
\index{NACA-Score}
\Z{\index{NACA}National Advisory Committee for Aeronautics

Diente ursprünglich der Bewertung von  Erkrankungen / Verletzungen und
der Entscheidung welche Soldaten zuerst auszufliegen wären. Heute
dient der NACA-Score der groben statistischen Erfassung.  
}
1 bis 7:


\begin{table}[H]
\Captionof{table}{NACA-Score.}

\begin{tabular}{cl}\addlinespace
\TabHeading{NACA}
&
\TabHeading{Bedeutung}
\\\addlinespace\midrule\addlinespace
 I & keine ärztliche Therapie erforderlich
\\\addlinespace
 II & ambulante Abklärung
\\\addlinespace
 III & stationäre Abklärung
\\\addlinespace
 IV & akute Lebensgefahr nicht ausgeschlossen
\\\addlinespace
 V & akute Lebensgefahr
\\\addlinespace
 VI & Reanimation
\\\addlinespace
 VII & Tod
\\\addlinespace\bottomrule
\end{tabular}
\end{table}
}{}{}{}{}



\clearpage
%%%%%%%%%%%%%%%%%%%%%%%%%%%%%%%%%%%%%%%%%%%%%%%%%%%%%%%%%%%
%%%%%%%%%%%%%%%%%%%%%%%%%%%%%%%%%%%%%%%%%%%%%%%%%%%%%%%%%%%
%%%%%%%%%%%%%%%%%%%%%%%%%%%%%%%%%%%%%%%%%%%%%%%%%%%%%%%%%%%
%%%%%%%%%%%%%%%%%%%%%%%%%%%%%%%%%%%%%%%%%%%%%%%%%%%%%%%%%%%
\section{Bewusstsein}
\index{Bewusstsein}

% \Definition{Überbegriff u.a. für Wachheit,
% Orientierung, Aufmerksamkeit, Auffassungsgabe, Denkverlauf und
% Merkfähigkeit\cite{Pschy259}.} 
% 
% \Z{Die "`wahre"' Definition
% von Bewusstsein ist seit tausenden von Jahren ungeklärt. Auch an dieser Stelle
% wird das Geheimnis nicht gelüftet werden, daher muss die obige Definition für
% den Alltag ausreichen.\footnote{Hinweisen zufolge könnte die Antwort eventuell
% "`42"' sein.}}

\protect\DefinitionExp{Bewusstsein}

bla

\protect\glsentrydesc{Bewusstsein}

Das Bewusstsein ist besonders wichtig hinsichtlich des Schutzes des
Menschen vor Bedrohung. Ein bewusstseinsklarer Mensch kann sich gegen
Gefahren wehren, ein eingetrübter oder bewusstloser Mensch kann dies
schlecht oder gar nicht. 

%%%%%%%%%%%%%%%%%%%%%%%%%%%%%%%%%%%%%%%%%%%%%%%%%%%%%%%%%%%
%%%%%%%%%%%%%%%%%%%%%%%%%%%%%%%%%%%%%%%%%%%%%%%%%%%%%%%%%%%
%%%%%%%%%%%%%%%%%%%%%%%%%%%%%%%%%%%%%%%%%%%%%%%%%%%%%%%%%%%
\subsection{Bewusstseinsstörungen}
\label{topic:bewusstseinsstoerung}

\Definition{Störung des Bewusstseins, besonders der
Wachheit und Orientierung.}


%%%%%%%%%%%%%%%%%%%%%%%%%%%%%%%%%%%%%%%%%%%%%%%%%%%%%%%%%%%
\Layout{0}{Allgemein}
{
Wie bereits oben erwähnt ist das Bewusstsein wesentlich an der Abwehr
von Gefahren beteiligt. Je mehr das Bewusstsein gestört oder
reduziert ist desto gefährlicher ist das für den Patienten. 

Die \Es{Zunge} kann zurück fallen und die Atemwege
verlegen. Dabei ist manchmal ein dem
\E{Schnarchen ähnliches} Atemgeräusch zu
vernehmen. 

Durch den Ausfall \index{Schutzreflexe!Ausfall} der
Schutzreflexe kann es zu einer
\Es{Aspiration}\index{Aspiration} von Speichel, Blut oder Erbrochenen kommen. 
}{
Gefahren:
\begin{itemize}
  \item Zurückfallen der Zunge.
  \item Ausfall der Schutzreflexe
  \item Aspiration.
\end{itemize}
}{}{}{}

%%%%%%%%%%%%%%%%%%%%%%%%%%%%%%%%%%%%%%%%%%%%%%%%%%%%%%%%%%%
\Layout{2}{Schutzreflexe}
{
Ab einem gewissen Punkt können auch so wesentliche
Funktionen wie der \Es{Schluck-}, \Es{Würge-} bzw. \Es{Hustenreflex}
ausfallen, und es kann Mageninhalt oder Blut in die
Lunge gelangen (\T{Aspiration}). Man spricht auch vom
\Es{Ausfall der Schutzreflexe}.

Bei schweren Bewusstseinsstörungen kann es auch zum Erschlaffen
wichtiger Muskeln wie z. B. der Zunge kommen. Fällt
diese zurück kann sie den Atemweg blockieren oder ganz verschließen.
}{}{}{}{}


\Fazit{Daher sind schwere Bewusstseinsstörungen
\Es{lebensgefährlich!}}


%%%%%%%%%%%%%%%%%%%%%%%%%%%%%%%%%%%%%%%%%%%%%%%%%%%%%%%%%%%
\Layout{0}{Einteilung der Störungen}
{
Man unterscheidet zwischen der Quantität ("`\E{Wieviel} Bewusstsein ist
vorhanden?"') und der Qualität ("`\E{Wie gut} funktioniert das vorhandene
Bewusstsein?"'). Maßgeblich ist der \Es{Bewusstseinsgrad} (quanititativ) und
die \Es{Orientiertheit} des Patienten (qualitativ). 

Die Glasgow Coma Scale ist ein spezielles Maßsystem für vorwiegend statistische
Zwecke. 
}{
\begin{itemize}
  \item Bewusstseinsgrad (quantitativ)
  \item Orientiertheit (qualitativ)
\end{itemize}
}{}{}{}

%%%%%%%%%%%%%%%%%%%%%%%%%%%%%%%%%%%%%%%%%%%%%%%%%%%%%%%%%%%
\Layout{2}{Bewusstseinsgrad}
{
\index{Bewusstseinsgrad}
\label{topic:bewusstseinsstoerung-grade}
\index{Bewusstsein!Symptome}

Der Bewusstseinsgrad gibt Auskunft über die Quantität des Bewusstseins ("`\E{Wieviel} Bewusstsein ist
vorhanden?"'). Er muß bei jedem Patienten festegestellt werden.

Die folgende Tabelle (\prettyref{tab:bewusstseinsgrade}) listet die möglichen
Bewusstseinsgrade auf:

\noindent\begin{minipage}{\linewidth}
\Captionof{table}{Bewusstseinsgrade\label{tab:bewusstseinsgrade}}

\begin{tabulary}{\linewidth}{LLLL}
\addlinespace
\TabSub{Bewusstseinsklar}
 &
 &
 wach
 &
\\\addlinespace
\TabSub{Bewusstseinsgetrübt}
&
\noindent\T{somnolent} \index{somnolent}
 &
schläfrig, aber leicht erweckbar.
 &
\\
 &
\noindent\T{soporös} \index{soporös}
 &
 kaum und nur mit erheblichen Aufwand (Schmerzreiz)
 erweckbar 
 &
\noindent\Es{Gefahr!}
\\\addlinespace
\TabSub{Bewußtlos}
 &
\noindent\T{komatös}\index{Koma}
 &
 nicht erweckbar
 &
\noindent\Ca{Lebensgefahr!}
\\\addlinespace\bottomrule
\end{tabulary}
\end{minipage}

}{
\begin{itemize}
  \item Klar.
  \item Getrübt.
  \begin{itemize}
    \item Somnolent.
    \item Soporös.
  \end{itemize}
  \item Bewusstlos -- Koma.
\end{itemize}

}{}{}{}

%%%%%%%%%%%%%%%%%%%%%%%%%%%%%%%%%%%%%%%%%%%%%%%%%%%%%%%%%%%
\Layout{2}{Orientierung}
{
\index{Orientierung}

Es ist wichtig zu Überprüfen ob der Patient orientiert ist.
Man unterscheidet dabei die Orientierung zu:  

\begin{labeling}[~--]{xxxxxxxxxxxxx}
  \item[Person] Name, Alter.
  \item[Zeit] Orientierung: Tageszeit,
  Wochentag, Monat, Jahreszeit, Jahr?
  \item[Ort] Wo befindet sich der Patient gerade (Stadt, Land, Adresse, Wohnung, Auto, etc.)?
  \item[Situation] Was passiert gerade?
\end{labeling}

}{
\begin{itemize}
    \item Zur Person. 
    \item Zeitlich. 
    \item Örtlich. 
    \item Zur Situation. 
\end{itemize}
}{}{}{}

%%%%%%%%%%%%%%%%%%%%%%%%%%%%%%%%%%%%%%%%%%%%%%%%%%%%%%%%%%%
\Layout{2}{Glasgow Coma Scale (GCS)}
{
\index{Glasgow Coma Scale}
\index{GCS}
\label{topic:gcs}

Die Glasgow Coma Scale ist eine spezielle Skala, welche der 'objektiven'
Einschätzung der Bewußtseinslage  dient. Sie hat einen Wertebereich von mind.~3
bis max.~15 Punkten. Bewertet werden die Augenöffnung, die beste verbale und
die beste motorische Reaktion des Patienten. 

Ein voll bewusstseinsklarer Patient hat 15 Punkte, ein schwerst
bewusstseinsgestörter oder toter Patient 3 Punkte (nicht 0!).

Heutzutage dient die \Abk{GCS} vorwiegend statistischen
Zwecken\ZFootnote{Regelmäßig findet man Handlungsanweisungen -- überwiegend
für ärztliches Personal -- in der Literatur, welche die \Abk{GCS} als
Entscheidungskriterium (z.\,B. für die Intubation) heranziehen. Der Nutzen ist
aber genauso regelmäßig umstritten.}.
}{
3\,--\,15 Punkte.
\begin{itemize}
    \item Augenöffnung.
    \item Beste verbale Reaktion.
    \item Beste motorische Reaktion.
\end{itemize}
Tafel \ref{tab:gcs}.
}{}{}{}


\begin{table}[H]
\Captionof{table}{GCS - Glasgow Coma Scale.\label{tab:gcs}}

\begin{tabulary}{\linewidth}{LCCL}
~ & \TabHeading{Max.} & \TabHeading{Punkte} &
\TabHeading{Reaktion}
\\\midrule
\TabSub{Augenöffnung} & 4 & 4 & spontan
\\
 & & 3 & auf Ansprache
\\
 & & 2 & Schmerzreiz
\\
 & & 1 & nicht
\\\midrule
\TabSub{Verbale Antwort} & 5 & 5 & klar
\\
 & & 4 & verwaschen
\\
 & & 3 & stammelnd
\\
 & & 2 & unartikuliert
\\
 & & 1 & nicht
\\\midrule
\TabSub{Motorische Reaktion} & 6 & 6 & spontan
\\
 & & 5 & gezielte Abwehrbewegungen
\\
 & & 4 & ungezielte Abwehrbewegungen
\\
 & & 3 & Beugung
\\
 & & 2 & Strecken
\\
 & & 1  & bewegt gar nicht
\\\bottomrule
\end{tabulary}
\end{table}




%%%%%%%%%%%%%%%%%%%%%%%%%%%%%%%%%%%%%%%%%%%%%%%%%%%%%%%%%%%
\Layout{0}{Ursachen}
{
Durch die Schädigung des zentralen Nervensystems kommt es zu der
Bewusstseinsstörung. Diese Schädigung kann direkt das \Abk{ZNS} betreffen
(primäre Ursache) oder über Umwege einen schädlichen Einfluß auf das \Abk{ZNS}
haben (sekundäre Ursache). \prettyref{tab:bewusstseinsstoerungen-ursachen}
listet typische Ursachen auf.

\begin{table}[H]
\Captionof{table}{Typische Ursachen von
Bewusstseinsstörungen.\label{tab:bewusstseinsstoerungen-ursachen}}

\begin{tabularx}{\linewidth}[H]{XX}
\TabHeading{Primäre Ursachen} & \TabHeading{Sekundäre Ursachen}
\\
{\tiny betreffen direkt das ZNS} & {\tiny betreffen über
Umwege das ZNS} 
\\\midrule
\begin{itemize}
    \item Verletzungen (SHT)
    \item Blutungen, Hirnschwellung
    \item Schlaganfall
    \item epileptischer Anfall
    \item Entzündungen
    \item \Z{Tumor}
    \item \ldots
\end{itemize}
 &
 \begin{itemize}
   \item Störungen der \ce{O2}-Versorgung:
   \begin{itemize}
     \item Atemstörungen  (Hypoxie)
     \item Herz-Kreislaufstörungen
   \end{itemize}
   \item Störungen des Wasserhaushaltes und
   Elektrolytentgleisungen
   \item Stoffwechselstörungen
   \begin{itemize}
     \item Zuckerhaushalt
     \item Leberversagen
   \end{itemize}
   \item Vergiftungen:
   \begin{itemize}
     \item Alkohol
     \item Medikamente, Drogen
     \item Reizgase,  Lösungsmittel ...
   \end{itemize}
   \item \ldots
 \end{itemize}
\\\bottomrule
\end{tabularx}
\end{table}
}
{
\begin{itemize}
  \item Direkt das \Abk{ZNS} betreffend (primär).
  \item Indirekt das \Abk{ZNS} betreffend (sekundär).
\end{itemize}

\prettyref{tab:bewusstseinsstoerungen-ursachen}.
}{}{}{}











\MassnahmenNG{Spezielle Maßnahmen}{Bewusstseinseintrübung}
{}{ 
% \paragraph{Maßnahmen
% bei getrübten Patienten}~

\begin{itemize}
  \item BAK-Schema  (Bewusstsein/Atmung/Kreislauf)
  \item Lagerung: situationsgerecht je nach
  Verdachtsdiagnose, im Zweifel \Es{stabile
  Seitenlage}, bei hochschwangeren Frauen auf die
  \E{linke} Seite.
  \item \TxBeiVit \TxMassVitBes
  
  \begin{itemize}
    \item \ce{O2}: nach Verdachtsdiagnose, im Zweifel
    ca. 6-8\LMin
  \end{itemize}
  
  \item Engmaschiges \Es{Monitoring} (Patient
  beobachten und regelmäßig ansprechen, Atmung,
  Pulsoxymetrie, RR)
  \item Wärmeerhalt
  \item \Es{Ursachenforschung} -- Warum ist der
  Patient eingetrübt?
  \item \Es{Neurocheck}
  \begin{itemize}
    \item Immer auch an \Es{BZ} denken!
  \end{itemize}
  \item \Es{Traumacheck}. Suche nach Primär- und
  Folgeverletzungen!
\end{itemize}

}

\MassnahmenNG{Spezielle Maßnahmen}{Bewusstlose und
soporöse Patienten}{}
{
\begin{itemize}
  \item BAK-Schema  (Bewusstsein/Atmung/Kreislauf)
  \item \TxMassVitBes
  \begin{itemize}
    \item Lagerung: stabile Seitenlage, bei hochschwangeren Frauen auf die
    \E{linke} Seite.
    \item \ce{O2}:  10\LMin , bei
    Zyanose\footnote{Zyanose: Bläuliche Verfärbung der Haut infolge einer unzureichenden
    Sauerstoffversorgung. Besonders auffällig an den Lippen, Gesicht, und
    Extremitäten.}  15\LMin
  \end{itemize}
  \item Wärmeerhalt
  \item \Es{Ursachenforschung} -- Warum ist der Patient bewusstlos oder
  soporös?
  \item \Es{Neurocheck}
  \begin{itemize}
    \item Immer auch an \Es{BZ} denken!
  \end{itemize}
  \item \Es{Traumacheck}. Suche nach Primär- und
  Folgeverletzungen!
\end{itemize}
}





\clearpage
%%%%%%%%%%%%%%%%%%%%%%%%%%%%%%%%%%%%%%%%%%%%%%%%%%%%%%%%%%%
%%%%%%%%%%%%%%%%%%%%%%%%%%%%%%%%%%%%%%%%%%%%%%%%%%%%%%%%%%%
%%%%%%%%%%%%%%%%%%%%%%%%%%%%%%%%%%%%%%%%%%%%%%%%%%%%%%%%%%%
%%%%%%%%%%%%%%%%%%%%%%%%%%%%%%%%%%%%%%%%%%%%%%%%%%%%%%%%%%%
\section{Atmung}

%%%%%%%%%%%%%%%%%%%%%%%%%%%%%%%%%%%%%%%%%%%%%%%%%%%%%%%%%%%
%%%%%%%%%%%%%%%%%%%%%%%%%%%%%%%%%%%%%%%%%%%%%%%%%%%%%%%%%%%
%%%%%%%%%%%%%%%%%%%%%%%%%%%%%%%%%%%%%%%%%%%%%%%%%%%%%%%%%%%
\subsection{Atemfrequenz, -tiefe und -volumen}
\label{topic:atemvolumina}


% STEU 2009-06-28: Trennung Oxygenation und Ventilation
\A{STEU 2009-06-28: Trennung Oxygenation und Ventilation}

\Layout{0}{Kennzahlen}
{
\index{Atemfrequenz}
\index{AF}
\index{Atemzugsvolumen}
\index{AZV}
\index{Atemminutenvolumen}  
\index{AMV} 

Es gibt 3 elementare Kennzahlen, welche die wichtigsten Atmungsparameter
beschreiben:

\begin{labeling}[\dots{}]{Atemminutenwolumen (AMV)}
  \item[\T{Atemfrequenz} (\T{AF})] gibt die Anzahl der Atemzüge pro
  Minute an.
  \item[\T{Atemzugsvolumen} (\T{AZV})] gibt die Menge
  Luft an, die \E{pro Atemzug} eingeatmet wird.
  \item[\T{Atemminutenvolumen} (\T{AMV})] gibt die Menge an \E{pro Minute}
  eingeatmeter Luft an.
\end{labeling}

Das Atemminutenvolumen kann aus der Atemfrequenz und dem Atemzugsvolumen
errechnet werden: 

\begin{equation}
\mbox{AMV} = \mbox{AZV} \times \mbox{AF}
\end{equation}

% Die
% \T{Atemfrequenz} (\T{AF}) gibt die Anzahl der Atemzüge pro
% Minute an. Das \T{Atemzugsvolumen} (\T{AZV}) gibt die Menge
% an Luft an die eingeatmet wird. Das \T{Atemminutenvolumen}
% (\T{AMV}) errechnet sich aus der Atemfrequenz und dem
% Atemzugsvolumen. Es gibt die Menge an \E{pro Minute}
% eingeatmeter Luft an.
}{
\begin{description}
    \item[AF] {\dotfill} Atemfrequenz
    \item[AZV] {\dotfill} Atemzugsvolumen
    \item[AMV] {\dotfill} Atemminutenvolumen
\end{description}

$\mbox{AMV} = \mbox{AZV} \times \mbox{AF}$
}{}{}{}



\Layout{0}{Normalwerte Atmung}
{
Die normale Atemfrequenz (\T{AF}) beträgt beim Erwachsenen
in Ruhe 12-16\,\nicefrac{x}{min}, das Atemzugsvolumen
(\T{AZV}) beträgt ca. 500\,ml. Das
entspricht einem Atemminutenvolumen (\T{AMV}) von  6\,--\,8\LMin.  Für eine
ausreichende Atmung ist auch eine ausreichende \ce{O2}-Konzentration der
Umgebungsluft notwendig! 
}{
\begin{description}
    \item[AF] {\dotfill} \Es{12-16\,\nicefrac{x}{min}}
    \item[AZV] {\dotfill} \Es{500\,ml}
    \item[AMV] {\dotfill} \Es{6-8\,l}
\end{description}
}{}{}{}

Auch wichtig ist die ausreichende
\ce{O2}-Konzentration in der Umgebungsluft. Ohne Sauerstoff in der Einatemluft
ist die beste Atmung nutzlos. Der normale \ce{O2}-Gehalt der Umgebungsluft
beträgt 21\,\%.
\marginlabel{\cite{KlinkePhysio3,SchmidtPhysio28,DRAn3,StriebelAn7}}

Die genauen Werte sind abhängig von den Faktoren Alter,
Geschlecht und Größe. 


\Fazit{AMV = AZV $\ast$ AF, beim Erwachsenen  ca. 6-8\,l}



\Co{Die Einstellungen eines Beatmungsgeräts können von den
Normalwerten abweichen, da man z.\,B. manchmal
"`vorsichtiger"', bzw. Lungen-schonender, beatmen möchte.}

Die Atemfrequenz ist stark altersabhängig:



\begin{table}[H]
\Captionof{table}{Atemfrequenz im Altersvergleich,
Richtwerte.}

\begin{tabular}{lrr}\addlinespace
 & \TabHeading{AF} & \TabHeading{AZV}
\\\addlinespace\midrule\addlinespace
\TabSub{Neugeborenes} & 40\,\nicefrac{x}{min} & 20-40\,ml 
\\\addlinespace
\TabSub{Kleinkind} & 25\,\nicefrac{x}{min} & 100-200\,ml
\\\addlinespace
\TabSub{Schulkind} & 20\,\nicefrac{x}{min} & 200-400\,ml
\\\addlinespace
\TabSub{Erwachsener} & 12\,--\,16\,\nicefrac{x}{min}  & ca. 500\,ml
\\\addlinespace\bottomrule
\end{tabular}
\end{table}

\Layout{0}{Totraum}
{
\index{Totraum}Unter dem Totraum
versteht man den Teil der Atemwege der nicht am
Gasaustausch teilnimmt. Dies ist im wesentlichen der gesamte Weg von der
Nasenspitze bis kurz vor den Lungenbläschen (Alveolen). Die
Luft muss diesen Totraum passieren um in die Lungenbläschen
zu gelangen! Der Totraum umfasst beim Erwachsenen ca.
150\,ml.\cite{KlinkePhysio3}

\index{Totraumventilation}\T{Totraumventilation}: Bei der
Totraumventilation ist das Atemzugsvolumen so klein dass
die "`frische"' Luft im Totraum bleibt und nicht in die
Lungenbläschen gelangt: Es kommt keine frische Luft in die
Lungenbläschen und folglich kein Sauerstoff ins Blut.

Bei der \index{Schnappatmung}\T{Schnappatmung} kann dies
passieren: Auf den ersten Blick sieht es so aus als würde
der Patient "`schnappend"' atmen, in Wirklichkeit ist die
Atmung so gering dass es nur zu einer Totraumventilation
kommt, d.h. Die Atmung ist nicht effektiv. 
}{
\begin{itemize}
    \item Totraum: Teil der Atemwege, der nicht am Gassaustausch teilnimmt.
    \item Totraumventilation.
    \item Schnappatmung = nicht ausreichend.
\end{itemize}
}{}{}{}




\Cave{Schnappatmung = keine Atmung}




\Layout{0}{Steuerung der Atmung}
{
Die Atmung wird durch das \Es{Atemzentrum im Hirnstamm} gesteuert. Daneben gibt
es Sensoren im Körper die den Gasgehalt des Blutes beobachten. 

Beim gesunden Patienten ist der \Es{Atemantrieb} vom \Es{\ce{CO2}-Gehalt} des
Blutes abhängig. Der Sauerstoffanteil spielt nur eine untergeordnete Rolle!
Unter \prettyref{topic:saeure-basen-haushalt} werden wir sehen dass das \ce{CO2}
eine wichtige Rolle im Säure-Basen-Haushalt des Körpers einnimmt.  

Wenn bei \Es{bestimmten Krankheitsbildern} die \E{\ce{CO2}-Konzentration
ständig hoch} ist, \E{gewöhnt} sich der Körper an diesen Zustand. Er beginnt
dann den \E{Atemantrieb direkt über den \ce{O2}}-Gehalt im Blut zu regeln.

Dies kann \Ca{Probleme} machen: Wird bei so einem Patienten der \ce{O2}-Gehalt
plötzlich künstlich erhöht fehlt der Atemantrieb und der Patient kann
\Es{Bewusstseinsstörungen} ("`\ce{CO2}-Narkose"') und einen
\Es{Atemstillstand} erleiden.

Das wichtigste Krankheitsbild bei dem so etwa vorkommen kann ist die
\Abk{COPD} (\ref{topic:copd}). 
}{
\begin{itemize}
  \item Atemzentrum im Hirnstamm
  \item Atemantrieb normalerweise durch \Es{\ce{CO2}}.
  \item Krankhaft:
  \begin{itemize}
    \item Wenn \ce{CO2}-Konzentration dauerhaft erhöht → Regelung durch
    \ce{O2}-Gehalt.
    \item Problem bei Sauerstoffgabe.
    \item V.a. bei \Abk{COPD} (\ref{topic:copd}).
  \end{itemize}
\end{itemize}
}{}{}{}






%%%%%%%%%%%%%%%%%%%%%%%%%%%%%%%%%%%%%%%%%%%%%%%%%%%%%%%%%%%
%%%%%%%%%%%%%%%%%%%%%%%%%%%%%%%%%%%%%%%%%%%%%%%%%%%%%%%%%%%
%%%%%%%%%%%%%%%%%%%%%%%%%%%%%%%%%%%%%%%%%%%%%%%%%%%%%%%%%%%
\subsection{Störungen der Atmung}

\Layout{0}{Einleitung}
{
Störungen der Atmung können aus vielen Ursachen eintreten. Tafel
\ref{tab:atemstoerungen-ursachen} beschreibt die wichtigsten Pathomechanismen
und Störungen. 

Der Begriff
\T{Ateminsuffizienz}\index{Ateminsuffizienz}\index{Insuffizienz!Atem-}
beschreibt allgemein eine \Es{nicht ausreichende Atemleistung}, unabhängig von der Ursache der Störung. Dies führt zu einem \ce{O2}-Mangel (Hypoxie). Oft hat der Patient das Gefühl nicht genug
Luft zu bekommen und leidet an \Es{Atemnot} (\T{Dyspnoe}\index{Dyspnoe}).

Die \T{Apnoe} (\Es{Atemstillstand}) ist die maximalste Form der
Ateminsuffizenz.
}{
\begin{itemize}
  \item Ateminsuffizienz \ldots{} nicht ausreichende Atmung
  \item Hypoxie \ldots{} \ce{O2}-Mangel
  \item Dyspnoe (Atemnot)
  \item Apnoe (Atemstillstand)
\end{itemize}
}{}{}{}


\Layout{0}{Gefahr der Ateminsuffizienz}
{
Bei einer Ateminsuffizienz sinkt der \Es{\ce{O2}-Gehalt} im Blut. Es kommt zu
einer \Es{Hypoxie} und in weiterer Folge zu einer Unterversorgung
lebenswichtiger Organe wie Herz und Hirn. 

Doch auch die \Es{\ce{CO2}-Abgabe} ist gestört und es kommt zu einer Anhäufung
von \ce{CO2} im Blut. In weiterer Folge kann es zu einer \E{atmungsbedingte
Übersäuerung} (\Es{Respiratorische Azidose}, 
\prettyref{topic:azidose-respiratorische}) kommen.
}{
\begin{itemize}
    \item \ce{O2}-Gehalt im Blut ↓ (Hypoxie).
    \item \ce{CO2}-Gehalt im Blut ↑.
    \item Respiratorische Azidose (\prettyref{topic:azidose-respiratorische}).
\end{itemize}
}{}{}{}

\AuthNotes{GABG: pathologische Atmungstypen inkl.
Zeichnung einbinden? Kussmaul, etc.}



\Layout{0}{Ursachen}
{
Tafel \ref{tab:atemstoerungen-ursachen}.
}{Tafel \ref{tab:atemstoerungen-ursachen}.}{}{}{}


\begin{table}[H]
\begin{gWideParEnv}
\Captionof{table}{Ursachen für
Atemstörungen.\label{tab:atemstoerungen-ursachen}}

\begin{tabulary}{\linewidth}[H]{LL}
\TabHeading{Störung des/der}  & \TabHeading{Vorkommen und Erkrankungen}
\\\addlinespace\midrule\addlinespace
\TabSub{\ce{O2}-Angebotes}
 &
Gärkeller (\ce{CO2}\,${\uparrow}$, \ce{O2}\,${\downarrow}$)

"`dünne Luft"' (z.\,B. Hochgebirge,

p\ce{O2}\,${\downarrow}$)

Defekte Gastherme (\ce{CO}\,${\uparrow}$,
\ce{O2}\,${\downarrow}$)

Ertrinken
\\\addlinespace
\TabSub{\ce{O2}-Diffusion}
 &
Lungenödem

Lungenentzündung

Lungenembolie
\\\addlinespace
\TabSub{Atemmechanik}
 &
Verlegung der oberen Atemwege (Zunge, Erbrochenes, Laryngospasmus,
         Glottisödem, Bolus)

Verlegung der unteren Atemwege (Bronchitis, Asthma,
Lungenödem)

Verminderung der Dehnbarkeit der Thoraxwand oder des Lungengewebes
     (Rippenfraktur, Pneumothorax, Pleuraerguß, Emphysem)
\\\addlinespace
\TabSub{Neuro\-muskulären Atemregulation}
 &
SHT (Schädel Hirn Trauma)

Schlaganfall (Insult)

Vergiftungen

Tumore

Hoher Querschnitt (Rückenmark durchtrennt)

Muskelerkrankungen

Nervenerkrankung

Stoffwechselstörungen
\\\addlinespace\bottomrule
\end{tabulary}
\end{gWideParEnv}
\end{table}


\Layout{0}{Symptome}
{
Atemstörungen können viele verschiedene Symptome verursachen. Tafel
\ref{tab:atmung-symptome} listet verschiedene Symptome auf. Neben technischen
Geräten ist vor allem \E{sehen}, \E{hören}, \E{fühlen} und ein
\E{aufmerksamer Blick} gefordert.
}{Tafel \ref{tab:atmung-symptome}.}{}{}{}


\begin{table}[H]
\Captionof{table}{Symptome: Atmung.\label{tab:atmung-symptome}}

\begin{tabulary}{\linewidth}[H]{LLLL}
\addlinespace
\TabHeading{Kriterium}
 &
\TabHeading{Befund}
 &
\TabHeading{Beschreibung}
 &
%\TabHeading{Anm.}
\\\addlinespace\midrule\addlinespace
\TabSub{Beschwerde}
 &
Atemnot (Dyspnoe) \index{Dyspnoe}
 &
Leitsymptom
 &
\\\addlinespace
\TabSub{Atemgeräusch}
 &
Brodeln \index{Atemgeräusch!brodelndes}
 &
Blubbern, klassisch für Lungenödem
\index{Atemgeräusch!blubberndes} &
\\\addlinespace
 &
Stridor \index{Stridor}
 &
 &
\\\addlinespace
 &
Brummen, Giemen
\index{Atemgeräusch!brummendes}\index{Atemgeräusch!giemendes}
& &
\\\addlinespace
 &
Rasselgeräusche \index{Rasselgeräusche}
 &
 &
\\\addlinespace
\TabSub{Frequenz}
 &
Beschleunigt
 &
Tachypnoe \index{Tachypnoe}
 &
\\\addlinespace
 &
Verlangsamt
 &
Bradypnoe \index{Bradypnoe}
 &
\\\addlinespace
\TabSub{Atemzugsvolumen}
 &
Schnappatmung
 &
 &
\\\addlinespace
 &
Flache Atmung
 &
 &
\\\addlinespace
 &
Tiefe Atmung
 &
 &
\\\addlinespace
\TabSub{Hautfarbe}
 &
Blass
 &
Normal, Anämie, Blutverlust
 &
\\\addlinespace
 &
Rosig
 &
Normal, \ce{CO}-Vergiftung
 &
\\\addlinespace
 &
Bläulich 
 &
Zyanose: Hypoxie!
 &
\\\addlinespace
\TabSub{Körperlich}
 &
Einziehungen an den Rippen
&
Einsatz der Atemhilfsmuskulatur
&
\\\addlinespace
&
Aufstützen
&
Einsatz der Atemhilfsmuskulatur
&
\\\addlinespace
&
Nasenflügeln \index{Nasenflügeln}
&
Atemnot, besonders bei Kleinkindern
&
\\\addlinespace
&
Aufrechte Position
&
&
\\\addlinespace
&
"`Paradoxe Atmung"'
&
Brustkorb senkt sich bei Einatmung
&
\\\addlinespace
\TabSub{Geräte}
&
Pulsoxymetrie
&
&
\\\addlinespace
&
\Z{Kapnometrie}
&
\Z{\ce{CO2}-Messung}
\\\addlinespace\bottomrule
\end{tabulary}
\end{table}




\MassnahmenNG{Spezielle Maßnahmen}{Atemstillstand}{}{
%\begin{itemize}
%    \item Freimachen der Atemwege
%    \begin{itemize}
%        \item Öffnen des Mundes \A{(Esmarch-Handgriff)
%        -- evtl. erklären}
%        \TempPicMissing{Esmarch-Handgriff}
%        \item Inspektion
%        \item fals nötig ausräumen
%        \item (falls nötig) Absaugen
%        \item Zungengrund wird angehoben -{\textgreater} Atemwege  sind frei!
%    \end{itemize}
%    \item Überprüfung der Atemtätigkeit (bei  überstreckter HWS)
%    durch
%    \begin{itemize}
%        \item SEHEN, HÖREN, FÜHLEN für mindestens 10  Sekunden!
%    \end{itemize}
%    \item Beatmung (keine ausreichende Spontanatmung bzw. bei Reanimation)
%    \begin{itemize}
%        \item Beatmungsbeutel mit Maske und
%        Reservoir (Erhöhung der Sauerstoffkonzentration)
%    \end{itemize}
%\end{itemize}
\Cave{Achtung: Bei einem mit Atemstillstand vorgefundenen
Patienten ist grundsätzlich von einer Reanimation
auszugehen.}

Vgl. die entsprechenden Kapitel im Teil Erste Hilfe und
Sanitätstechnik bezüglich der genauen Durchführung der
Maßnahmen. 

\Z{Nur in sehr seltenen Fällen kann primär ein isolierter
Atemstillstand ohne gleichzeitigem Kreislaufstillstand
beobachtet werden. Auf diese Spezialfälle wird hier nicht
weiter eingegangen. Auch der iatrogene Atemstillstand (z.B
im Rahmen einer Narkose) wird an dieser Stelle nicht
behandelt.} }








\clearpage
%%%%%%%%%%%%%%%%%%%%%%%%%%%%%%%%%%%%%%%%%%%%%%%%%%%%%%%%%%%
%%%%%%%%%%%%%%%%%%%%%%%%%%%%%%%%%%%%%%%%%%%%%%%%%%%%%%%%%%%
%%%%%%%%%%%%%%%%%%%%%%%%%%%%%%%%%%%%%%%%%%%%%%%%%%%%%%%%%%%
%%%%%%%%%%%%%%%%%%%%%%%%%%%%%%%%%%%%%%%%%%%%%%%%%%%%%%%%%%%
\section[Herz- / Kreislauf]{Herz-\,/\,Kreislauf}
\label{bkm:RefHeading42514202}

\begin{minipage}{\linewidth}
\Captionof{table}{Gegenüberstellung der
wichtigsten Kennzahlen, die die Herz- und Lungenleistung
beeinflussen.}
\begin{tabularx}{\linewidth}{XXX}
\TabHeading{~} & \TabHeading{Herz} & \TabHeading{Lunge}
\\\addlinespace\toprule\addlinespace
\TabSub{Was?} & Blut & Luft
\\\addlinespace
\TabSub{Qualität} & Sauerstoffgehalt\Z{\footnote{\Z{Nicht
gleichzusetzen mit der Sauerstoffsättigung}}} &
Sauerstoffkonzentration 
\\\addlinespace\midrule\addlinespace
\TabSub{Wie oft?}
& 
Herzfrequenz 

(HF) 
&
Atemfrequenz 

(AF)
\\\addlinespace
\TabSub{Wieviel?}

(pro Aktion)
&
Schlagvolumen

(SV)
&
Atemzugsvolumen

(AZV)
\\\addlinespace\midrule\addlinespace
\TabSub{Erzielte Menge}
&
Herzminutenvolumen 

(HMV)
&
Atemminutenvolumen

(AMV)
\\\addlinespace\bottomrule
\end{tabularx}
\end{minipage}

%%%%%%%%%%%%%%%%%%%%%%%%%%%%%%%%%%%%%%%%%%%%%%%%%%%%%%%%%%%
%%%%%%%%%%%%%%%%%%%%%%%%%%%%%%%%%%%%%%%%%%%%%%%%%%%%%%%%%%%
%%%%%%%%%%%%%%%%%%%%%%%%%%%%%%%%%%%%%%%%%%%%%%%%%%%%%%%%%%%
\subsection{Blutdruck} 
\label{topic:blutdruck}

Durch die Pumpfunktion des Herzens entsteht ein
lebenswichtiger Druck im (arteriellen) Gefäßsystem, der
\E{\Gls{blutdruck}}.\index{Blutdruck}

test:

\glsentrydesc{blutdruck}

:ende test

\Layout{0}{Blutdruck}
{
Durch die rhythmische Funktion des
Herzens kommt es infolge des Blutauswurfes zu genauso rhythmischen Druckveränderungen.
Während des Auswurfs von Blut in das
Gefäßsystem (Herzanspannung, Systole)
spricht man vom systolischen Blutdruck, während
der Entspannungsphase des Herzens (Diastole) vom
diastolischen Blutdruck. Natürlich ist der Blutdruck bei
der Anspannung des Herzens höher als bei der Entspannung
(\E{Systolischer Blutdruck > Diastolische Blutdruck)}.
}{
\begin{description}
    \item[systolisch:] Maximaldruck in der Arterie am Höhepunkt der
Kammerkontraktion
    \item[diastolisch:] Druck in der Arterie während die Kammer erschlafft ist
\end{description}
}{}{}{}


Zur Durchführung der Blutdruckmessung vgl.
\prettyref{topic:rr-messung}.



%%%%%%%%%%%%%%%%%%%%%%%%%%%%%%%%%%%%%%%%%%%%%%%%%%%%%%%%%%%
% TODO : Section: Layout-Anpassung
\TODO{GABS}{2010-mm-yy}{Section: Layout-Anpassung notwendig.}
   {Dieser Abschnitt entspricht nicht dem Standard-Layout!
}
%%%%%%%%%%%%%%%%%%%%%%%%%%%%%%%%%%%%%%%%%%%%%%%%%%%%%%%%%%%
\begin{description}
    \item[Normbereich] zwischen 90/60\,mmHg und
    140/90\,mmHg \index{Blutdruck!Normbereich}
    \item[Hypotonie] RR {\textless}\,90/60\,mmHg
    \index{Hypotonie}
    \item[Hypertonie] RR {\textgreater}\,140/90\,mmHg
    \index{Hypertonie}
\end{description}

{\tiny Werte beziehen sich auf den Erwachsenen.}

Achtung: Eine Hypertonie ist zwar eine Erkrankung, aber nicht immer im
Einsatz von Bedeutung. Viele Patienten sind an ihren erhöhten
Blutdruck gewohnt und haben keine Beschwerden. Es darf nie nur nach dem
Wert gehandelt werden, es muss auch immer bedacht werden: 

\begin{itemize}
    \item \Es{Zustand} des Patienten.
    \item Der \Es{sonst übliche} Blutdruck des Patienten
    (erfragen, Blutdruckprotokoll).
    \item Der \Es{Verlauf} des Blutdruckes während der
    Behandlung (Traumapatienten!).
    \item Der \Es{Erregungszustand} des Patienten.
\end{itemize}

\minisec{Beispiele}

\emph{Ein 65-jähriger, männlicher Patient klagt über
Schwindelgefühl, sie messen einen Blutdruck von 100/70. Im Blutdruckprotokoll finden sie
für die vergangenen vier Tage folgende Einträge: 170/110, 180/115,
175/113, 183/112, ... .}


Wie beurteilen Sie den Blutdruck?
\index{Blutdruck!Beurteilung}

\begin{quote}
Der Blutdruck ist für den Patienten zu niedrig, der Körper ist auf
eine so plötzliche Umstellung nicht vorbereitet gewesen. Im
Anamnesegespräch erfahren sie, dass er heute das erste Mal einen
"`Nitro-Spray"' verwendet hat. Sie
vermerken dies am Einsatzprotokoll. Im Krankenhaus erfahren Sie, dass
wahrscheinlich eine Überdosierung dieses Medikaments die Ursache
dieses Blutdruckabfalls war. ${\rightarrow}$ Obwohl der Patient laut
Lehrbuch einen "`schönen"' Blutdruck
hatte, war dieser dennoch zu niedrig!
\end{quote}

\emph{Ein 56-jähriger Patient klagt über starke
kolikartige Schmerzen im linken Unterbauch. Sie messen einen Blutdruck von 155/95. Der Patient
kennt seinen "`gewöhnlichen"' Blutdruck
nicht.}  

Wie beurteilen Sie den Blutdruck? 

\begin{quote}
Der Blutdruck ist laut Lehrbuch erhöht, aber dies ist erklärbar die
starken Schmerzen und den damit verbundenen Erregungszustand. Die
Beschwerden stehen wahrscheinlich in keinem ursächlichen Zusammenhang
mit den Blutdruck. 
\end{quote}

\emph{Eine 50-jährige Patientin hat Nasenbluten. Sie
messen einen Blutdruck von 220/120. Auf Nachfrage gibt sie an etwas Kopfsschmerzen zu haben,
und etwas schwindlig sei sie auch.} 

Wie beurteilen Sie den Blutdruck? 

\begin{quote}
Der Blutdruck ist sehr stark erhöht, außerdem zeigt die
Patientin Zeichen einer hypertensiven Krise (Nasenbluten, Kopfweh, Schwindel).
Hier ist der erhöhte Blutdruck Ursache für die Beschwerden und muss
mittels Medikamente unbedingt gesenkt werden. 
\end{quote}


%%%%%%%%%%%%%%%%%%%%%%%%%%%%%%%%%%%%%%%%%%%%%%%%%%%%%%%%%%%
%%%%%%%%%%%%%%%%%%%%%%%%%%%%%%%%%%%%%%%%%%%%%%%%%%%%%%%%%%%
%%%%%%%%%%%%%%%%%%%%%%%%%%%%%%%%%%%%%%%%%%%%%%%%%%%%%%%%%%%
\subsection{Kreislaufstillstand}\label{bkm:RefHeading42514602}\label{topic:kreislaufstillstand}
\index{Kreislaufstillstand}

\Definition{Totales Pumpversagen des Herzens.} 

%%%%%%%%%%%%%%%%%%%%%%%%%%%%%%%%%%%%%%%%%%%%%%%%%%%%%%%%%%%
% TODO : Überarbeitung
\TODO{GABS}{2010-mm-yy}{Überarbeitung notwendig.}
   {Dieser Teil ist nicht gut formuliert und/oder unvollständig, 
   er benötigt eine Überarbeitung!
}
%%%%%%%%%%%%%%%%%%%%%%%%%%%%%%%%%%%%%%%%%%%%%%%%%%%%%%%%%%%

Das Pumpversagen kann entstehen durch:

\begin{itemize}
    \item Asystolie (keine Herztätigkeit)
    \item Kammerflimmern (unkoordinierte Herztätigkeit)
    \item Zu schnelle Herzfrequenz, bei der sich das Herz
    nicht mehr füllen oder etwas auswerfen kann ((Pulslose) Ventrikuläre Tachykardie)
    \item Elektromechanische Entkopplung: keine
    Pumpreaktion auf Signale des Reizleitungssystems im Herz
\end{itemize}

\Layout{2}{Ursachen}
{
\begin{itemize}
  \item Minderversorgung des Herzmuskels mit \ce{O2}
  \begin{itemize}
    \item Herzinfarkt
    \item Herzinsuffizienz
  \end{itemize}
  \item Störungen der Atmung
  \begin{itemize}
    \item Asthma bronchiale
    \item Lungenembolie
    \item zentrale Atemregulationsstörungen
  \end{itemize}
  \item sonstige Störungen
  \begin{itemize}
    \item Trauma
    \item Vergiftungen
    \item thermische Schäden
    \item Elektrolytverschiebungen,...
  \end{itemize}
\end{itemize}
}{}{}{}{}

\Layout{2}{Folgen und Symptome}
{
%%%%%%%%%%%%%%%%%%%%%%%%%%%%%%%%%%%%%%%%%%%%%%%%%%%%%%%%%%%
% TODO : Fließtext/Slide fehlt!
\TODO{GABS}{2010-mm-yy}{Fließtext/Slide fehlt!}
   {Fließtext/Slide fehlt!}
%%%%%%%%%%%%%%%%%%%%%%%%%%%%%%%%%%%%%%%%%%%%%%%%%%%%%%%%%%%
\begin{itemize}
  \item Aussetzen des Herzschlages $\rightarrow$ Patient ist
  sofort pulslos!
  \item Bewusstlosigkeit (nach 15-20\,sek)
  \item Atemstillstand/Schnappatmung (nach ca. 10\,sek)
  \item \Es{Irreversible Hirnschäden nach 3-5\,min}, bei
  Kälte allerdings deutlich länger!
\end{itemize}
}{}{}{}{}


\A{KOCH: hab ich die (Hirnschäden) nicht aufgrund des
Atemstillstandes?

GABS: nein, auch ohne Atmung erfolgt noch eine Zeit lang
eine ausreichende \ce{O2}-Versorgung, die Ischämie beginnt
erst mit dem Kreislaufstillstand.}







%%%%%%%%%%%%%%%%%%%%%%%%%%%%%%%%%%%%%%%%%%%%%%%%%%%%%%%%%%%
%%%%%%%%%%%%%%%%%%%%%%%%%%%%%%%%%%%%%%%%%%%%%%%%%%%%%%%%%%%
%%%%%%%%%%%%%%%%%%%%%%%%%%%%%%%%%%%%%%%%%%%%%%%%%%%%%%%%%%%
\subsection{\Z{Herzinsuffizienz, Herzversagen}}
\index{Herzinsuffizienz}\index{Herzversagen}

Weitere, oft auch sehr bedrohliche, Störungen des Herzens werden im
Kapitel \emph{"'Erkrankungen und medizinische Notfälle"'}
unter \prettyref{topic:herz-kreislauf-stoerungen} besprochen. 




%%%%%%%%%%%%%%%%%%%%%%%%%%%%%%%%%%%%%%%%%%%%%%%%%%%%%%%%%%%
%%%%%%%%%%%%%%%%%%%%%%%%%%%%%%%%%%%%%%%%%%%%%%%%%%%%%%%%%%%
%%%%%%%%%%%%%%%%%%%%%%%%%%%%%%%%%%%%%%%%%%%%%%%%%%%%%%%%%%%
%%%%%%%%%%%%%%%%%%%%%%%%%%%%%%%%%%%%%%%%%%%%%%%%%%%%%%%%%%%
\section{Schock}
\index{Schock}
\label{topic:schock}

\Definition{"`Lebensbedrohliche Kreislaufstörung aufgrund von
Unterversorgung des Körpers mit Blut und Sauerstoff infolge des
Versagens des Kreislaufsystems."'}


\DefinitionExp{Schock}

\Fazit{Es können einzelne oder mehrere Teile des
Kreislaufsystems betroffen sein  (Blut, Gefäße, Herz).}

\Fazit{Ob ein Schock oder eine Schockgefahr besteht wird anhand von
\underline{Schocksymptomen} und anhand der
\underline{Verdachtsdiagnose} beurteilt.}

% TODO Graphik Schockentwicklung RR und HF

%%%%%%%%%%%%%%%%%%%%%%%%%%%%%%%%%%%%%%%%%%%%%%%%%%%%%%%%%%%
% TODO : Fließtext/Slide fehlt!
\TODO{GABS}{2010-mm-yy}{Fließtext/Slide fehlt!}{Fließtext/Slide fehlt!}
%%%%%%%%%%%%%%%%%%%%%%%%%%%%%%%%%%%%%%%%%%%%%%%%%%%%%%%%%%%
Schock bedeutet:

\begin{itemize}
  \item Akute Kreislaufinsuffizienz =
  \Es{Missverhältnis
  Herzauswurf\,:\,Sauerstoffbedarf} des Körpers
  \item Mögliche Ursachen:
  \begin{enumerate}
    \item Herz pumpt zu wenig
    \item Zu wenig Blut bzw. Flüssigkeit vorhanden
    \item Gefäße sind zu weit gestellt und Blut versackt
  \end{enumerate}
\end{itemize}

\Fazit{Im Endeffekt kommt immer zu wenig Blut mit
Sauerstoff zu den Organen. Dadurch ist die
Ge\-webedurchblutung für die Versorgung aller Organe mit
Sauer\-stoff nicht ausreichend.}

Körpereigene Mechanismen können versuchen gegenzusteuern
um den Schock zu kompensieren. Dies geht allerdings nur eine
gewisse Zeit und in einem gewissen Ausmaß gut, daher ist eine
schnell  einsetzende Therapie wichtig! Bei länger bestehendem Schock
kann dieser Zustand irreversibel werden. 

\E{Kein} Schock besteht bei kurzzeitigen 
Bewusstseinsstörungen wie Kollaps/Synkope
(\underline{kurzfristige} Blutverteilungsstörung,
kurzfristige Minderdurchblutung des Gehirns,  flüchtige Symptome)

\Layout{0}{Kompensation}
{
\index{Schock!Kompensation}
Der Körper versucht selber dagegen etwas zu tun und gegenzusteuern
("`kompensieren"'). Der
wichtigste Mechanismus ist die
\index{Zentralisation}\T{Zentralisation}. Darunter versteht
man die Blutumverteilung zugunsten lebenswichtiger Organe wie
Herz, Gehirn, Lunge und Leber. Dabei werden z.\,B. Haut, Muskel
und Darm
sowie die Extremitäten (=Peripherie) kaum mehr  durchblutet.
Periphere Venen kollabieren und die Rekapillarisierungszeit
verlängert sich. Die Hautdurchblutung wird verringert und
die Haut wird kühl und bleich.

Wenn die Kompensationsmaßnahmen 
nicht mehr ausreichen, spricht man vom \E{"`dekompensierter
Schock"'}\index{Schock!dekompensierter}. Es besteht
spätestens dann akute \Es{Lebensgefahr}!
}{
\begin{itemize}
    \item "`Gegensteuerung"'.
    \item Zentralisation.
    \item Ende: "`Dekompensierter Schock"'.
\end{itemize}
}{}{}{}



Beachte: 

\Cave{Bei jedem Patienten ist zu erheben, ob ein Schock
oder eine realistische Schockgefahr besteht oder nicht bzw. festzustellen um welche Art von
Schock es sich handelt (s.u.)!}

\Cave{Ein Schockzustand ist grundsätzlich als vitale
Bedrohung einzustufen!}

%%%%%%%%%%%%%%%%%%%%%%%%%%%%%%%%%%%%%%%%%%%%%%%%%%%%%%%%%%%
%%%%%%%%%%%%%%%%%%%%%%%%%%%%%%%%%%%%%%%%%%%%%%%%%%%%%%%%%%%
%%%%%%%%%%%%%%%%%%%%%%%%%%%%%%%%%%%%%%%%%%%%%%%%%%%%%%%%%%%
\subsection{Allgemeine Schocksymptomatik}
\label{topic:schock-allgemeine-symptomatik}
\index{Schock!allgemeine Symptomatik}

\begin{table}[H]
\Captionof{table}{Allgemeine Schocksymptomatik.}

\begin{tabular}{p{3.5cm}p{8cm}}
\addlinespace
\TabHeading{Stadium}
&
\TabHeading{Symptome}
\\\addlinespace\midrule\addlinespace
{\sffamily\Large\color{DarkGreen} \index{Schockgefahr}Schockgefahr}

\centering\includegraphics[width=1cm]{./icons/sm-basic.png}
&
 \begin{itemize}
     \item Keine eindeutigen Anzeichen
     \item Vorausschauend handeln!
     \item Verdachtsdiagnose bedenken.  Muss man damit
     rechnen, dass sich ein Schock entwickelt?  Besteht
     Schockgefahr?
 \end{itemize}
\\\addlinespace
{\sffamily\Large\color{OrangeRed}Schock}

\centering\includegraphics[width=1cm]{./icons/sm-pale.png}
&
 \begin{itemize}
     \item HF ${\uparrow}$
     \item mäßiger RR-Abfall, aber RR noch {\textgreater} HF
     \item neurologisch unauffällig, ev. euphorisch oder
     agitiert
     \item Haut blass, evtl. kaltschweißig
     \item Venen sind kollabiert oder gestaut
     \item Peripherie schlecht durchblutet, Rekapzeit erhöht
 \end{itemize}
\\\addlinespace
{\sffamily\Large\color{red}Schwerer Schock

Vollbild}

\centering\includegraphics[width=1cm]{./icons/sm-pale-frown.png}
&
\begin{itemize}
  \item Tachykardie (HF ${\uparrow}$, "`fliehender"' Puls)
  \item RR-Abfall (schwacher Puls)
  \begin{itemize}
    \item Herzfrequenz {\textgreater} systolischer RR
  \end{itemize}
  \item Tachypnoe (schnelle, flache Atmung)
  \item neurologisch auffällig, getrübt
  \item Haut blass, zyanotisch (bläulich), feucht
  \item Venen kollabiert oder gestaut
  \item nur mehr schwer behebbar! ${\rightarrow}$ Schock  möglichst
  frühzeitig erkennen \&  behandeln
\end{itemize}
\\\addlinespace
{\sffamily\Large\color{LightSlateGray}Endstadium}

\huge\centering $\skull$
&
\begin{itemize}
  \item Haut grau, klebrig
  \item Versagen der Herzfunktion
  \item RR kaum mehr messbar
  \item evtl. Koma
  \item evtl. Schnappatmung
  \item Patient ist dabei zu sterben.
\end{itemize}
\\\addlinespace\bottomrule
\end{tabular}
\end{table}

\Fazit{Das erste Anzeichen eines beginnenden Schocks ist
oft eine kühle Haut\footnote{Natürlich kann es auch andere
Ursachen für eine kühle Haut geben, z.\,B. eine kalte
Umgebungstemperatur}\cite{Kaplan2001}!}

\Layout{0}{Schockschere}
{
\index{Schock!-schere}
\label{topic:schockschere}
Normalerweise ist die Herzfrequenz niedriger als der systolische
Blutdruck. Wenn der Körper versucht den Schock zu
kompensieren, kann sich dieses Verhältnis umkehren. Man spricht dann von der
Schockschere, der Patient befindet sich dann \Es{bereits  im Vollbild
eines Schocks!} \E{Die Schockschere ist ein \Es{spätes Zeichen!} Auch wenn sich
ein Patient (noch)) nicht in der Schockschere befindet, kann er trotzdem einen
Schock haben!}

Achtung: \Es{Kinder haben andere Normalwerte für HF und
RR!} Die Faustregel ist deshalb nicht anwendbar!

% % TODO Schockschere
% \Captionof{figure}{Schockschere. \AuthNotes{Achtung! Bild
% fehlt!}}

\Fazit{Wenn sich das Verhältnis 'Herzfrequenz zu Systolischem
Blutdruck' umkehrt, befindet sich der Patient bereits in
einem schweren Schock.} 
\Fazit{Die Schockschere ist ein \Es{Spätzeichen!} Auch wenn die Herzfrequenz
höher ist als der systolische Blutdruck kann sich der Patient bereits im Schock
befinden!}
\Fazit{Ob ein Schock oder eine Schockgefahr besteht wird anhand von
\E{Schocksymptomen} und anhand der
\E{Verdachtsdiagnose} beurteilt, \Es{\E{nicht}} primär anhand der Schockschere!}
}{
\begin{itemize}
  \item RR\textsubscript{syst}\,{\textless}\,HF.
  \item Dann bereits \E{schwerer} Schock.
  \item Spätzeichen, auch ohne Schockschere kann Schock bestehen. 
  \item[\Ca{!}] Schockbeurteilung: Schocksymptome und Verdachtsdiagnose,
  \E{nicht} primär Schockschere.
  \item Für Kinder nicht anwendbar. 
\end{itemize}
}{}{}{}


\subsection{Allgemeine Maßnahmen der Schockbekämpfung}
\label{topic:schockbekaempfung}

\Fazit{Grundsatz: \Es{Der Schock muss bekämpft werden
bevor man ihn bemerkt.}}

\MassnahmenNG{Allgemeine Maßnahmen}{Schockbekämpfung}
{\label{m:am-schock}}
{
\begin{enumerate}
  \item \TxMassVitBes
  \begin{itemize}
    \item Situationsgerechte  Lagerung, je nach Diagnose und Schockart
    \item Sauerstoffgabe hochdosiert (ab 8\LMin{} bei Schockgefahr,
    15\LMin{} oder mehr bei bestehendem Schock, je nach Schwere)
    \item Genaues, ständiges Monitoring
  \end{itemize}
  \item Beengende Kleidung öffnen
  \item Wärmeerhalt
  \item Psychische Betreuung
  \item Weitere Maßnahmen je nach Schockart
\end{enumerate}
}
    \AuthNotes{Freigabe: Sauerstoffgabe hochdosiert (ab 6l/min bei
    Schockgefahr, 10l/min oder mehr bei bestehendem Schock, je nach Schwere}



%%%%%%%%%%%%%%%%%%%%%%%%%%%%%%%%%%%%%%%%%%%%%%%%%%%%%%%%%%%
%%%%%%%%%%%%%%%%%%%%%%%%%%%%%%%%%%%%%%%%%%%%%%%%%%%%%%%%%%%
%%%%%%%%%%%%%%%%%%%%%%%%%%%%%%%%%%%%%%%%%%%%%%%%%%%%%%%%%%%
\subsection{Schock -- Arten}
\index{Schock!Arten}

\begin{center}
\begin{tabularx}{\linewidth}[H]{XX}
\toprule
\Es{Mechanismus}
 &
\Es{Schockart}
\\\midrule
Herz pumpt zu wenig
&
{\begin{itemize}
    \item Kardiogener Schock
\end{itemize}}
\\
Zu wenig Flüssigkeit\,/\,Blut vorhanden 
 &
{\begin{itemize}
    \item Hypovolämischer Schock
\end{itemize}}
\\
Gefäße inadäquat weitgestellt, Blut versackt
 &
{\begin{itemize}
    \item Anaphylaktischer Schock
    \item Septischer Schock, toxischer Schock
    \item Neurogener Schock
\end{itemize}}
\\\bottomrule
\end{tabularx}
\end{center}


\subsubsection{Hypovolämischer Schock}
\label{topic:schock-hypovolaemischer}
\index{Schock!Hypovolämischer}

\Definition{Schock aufgrund eines Verlustes von Blutvolumen aus den
Gefäßen.}


%%%%%%%%%%%%%%%%%%%%%%%%%%%%%%%%%%%%%%%%%%%%%%%%%%%%%%%%%%%
% TODO : Fließtext/Slide fehlt!
\TODO{GABS}{2010-mm-yy}{Fließtext/Slide fehlt!}{Fließtext/Slide fehlt!}
%%%%%%%%%%%%%%%%%%%%%%%%%%%%%%%%%%%%%%%%%%%%%%%%%%%%%%%%%%%

Es kann dabei verloren gehen: 

\begin{itemize}
    \item Blut
    \item Plasma
    \item Wasser und Elektrolyte
\end{itemize}

\paragraph{Ursachen}
%%%%%%%%%%%%%%%%%%%%%%%%%%%%%%%%%%%%%%%%%%%%%%%%%%%%%%%%%%%
% TODO : Fließtext/Slide fehlt!
\TODO{GABS}{2010-mm-yy}{Fließtext/Slide fehlt!}{Fließtext/Slide fehlt!}
%%%%%%%%%%%%%%%%%%%%%%%%%%%%%%%%%%%%%%%%%%%%%%%%%%%%%%%%%%%



\begin{center}
\begin{tabularx}{\linewidth}[H]{XX}
\toprule
Flüssigkeitsverlust nach außen
&
Flüssigkeitsverlust nach innen
\\\toprule
\begin{itemize}
  \item Blutung
  \item gastrointestinal (Durchfall, Diarrhoe)
  \item Verbrennungen (Hautbarriere versagt  groß\-flächig)
  \item renal (Diabetes mellitus, Diuretika, Polyurie bei  Nierenversagen)
\end{itemize}
&
\begin{itemize}
  \item Innere Blutungen
  \begin{itemize}
    \item Nach Trauma, rupturierte Eileiter\-schwanger\-schaft, etc
  \end{itemize}
  \item Wassereinlagerung ins Gewebe
  \item \Z{Aszites
  ("`Wasserbauch"' bei
  Lebererkrankungen  etc.)}
\end{itemize}
\\\bottomrule
\end{tabularx}
\end{center}
Bei akutem Flüssigkeitsverlust kann ein Verlust von bis zu 20\,\% des
Blutvolumens  gut kompensiert werden. Auf größere Verluste  folgt
ein Absinken von Herzauswurf und RR!

\MassnahmenNG{Spezielle Maßnahmen}{Hypovolämischer Schock}{}
{

\begin{itemize}
  \item Allgemein: \Es{Allgemeine Maßnahmen der
  Schockbekämpfung} (s.
  \prettyref{m:am-schock})
  \item Blutstillung  (weitere Verluste verhindern!)
  \item ggfs. Trauma-Vorgehen beachten
  \item Wenn möglich schonende Bergung
  \item Lagerung: grundsätzlich Beine hoch
  
  Wichtige Ausnahmen: {\textperiodcentered}\,kardiale Ursache,
  {\textperiodcentered}\,Kopf- / Brust\-ver\-letzungen,
  {\textperiodcentered}\,Hirndruck (SHT),
  {\textperiodcentered}\,Knochenbrüche der Beine oder des
  Beckens,
  {\textperiodcentered}\,Wirbelsäulenverletzungen,
  {\textperiodcentered}\,Bewusstlosigkeit ${\rightarrow}$ sonst Flachlagerung oder stabile Seitenlage \A{KOCH: darüber müssen wir reden: Schock vor Schädel?}
  %TODO Schock vor Schädel
  \item evtl. Aviso (Ankündigung) im  Krankenhaus / Schockraum
  \item Abtransport: rasch, schonend, nicht überstürzt
\end{itemize}
}

\AuthNotes{Begriff "`Schocklagerung"'
ein für alle mal in der Ausbildung streichen, nur Hinweis dass dieser Ausdruck
veraltet ist sollten sie ihn wo aufschnappen. Stattdessen
"`Beine-hoch"'-Lagerung.

Wenn KI gegen Beine hoch: Flachlagerung bzw. stab. Seitenlage empfehlen?

STEU: ja, Flachlagerung wenn KI. 2009-06-28}


\subsubsection{Kardiogener Schock}
\index{Schock!Kardiogener}
\label{topic:schock-kardiogener}

\Definition{Schock infolge mangelnder Herzpumpleistung.}
(Die \E{Pulmonalembolie} als Auslöser eines Schocks wird
auch zum kardiogenen Schock gerechnet, da es durch die
Verstopfung des kleinen Kreislaufes zu einer mangelhaften
Füllung des Herzens und damit zu einer mangelhaften
Auswurfleistung kommt.)

%%%%%%%%%%%%%%%%%%%%%%%%%%%%%%%%%%%%%%%%%%%%%%%%%%%%%%%%%%%
% TODO : Überarbeitung
\TODO{GABS}{2010-mm-yy}{Überarbeitung notwendig.}
{Dieser Teil ist nicht gut formuliert und/oder unvollständig, 
er benötigt eine Überarbeitung!}
%%%%%%%%%%%%%%%%%%%%%%%%%%%%%%%%%%%%%%%%%%%%%%%%%%%%%%%%%%%

\begin{center}
\begin{tabularx}{\linewidth}[H]{XX}
\toprule
Ursachen &
charakterisiert durch\\\midrule
\begin{itemize}
  \item Herzversagen, Herzinsuffizienz
  \begin{itemize}
    \item Myokardinfarkt (häufigste Ursache)
    \item Herzrhythmusstörungen
  \end{itemize}
  \item Lungenembolie
\end{itemize}
 &
 \begin{itemize}
   \item RR-Abfall
   \item eingeschränkte Pumpleistung
   \begin{itemize}
     \item Blut staut sich zurück
   \end{itemize}
 \end{itemize}
\\\bottomrule
\end{tabularx}
\end{center}

\MassnahmenNG{Spezielle Maßnahmen}{Kardiogener Schock}{}{

\begin{itemize}
  \item Allgemein: \Es{Allgemeine Maßnahmen der
  Schockbekämpfung}
  (\prettyref{m:am-schock})
  \item \TxMassVitBes
  \begin{itemize}
    \item \ce{O2}: 10\LMin{}   oder mehr
    \item Lagerung: Oberkörper hoch! Evtl. Beine hängen  lassen.
  \end{itemize}
  \item Medikamentöse Therapie durch Notarzt
  \Z{(Diuretika, Katecholamine, Heparin)}
\end{itemize}
}







\subsubsection{Anaphylaktischer Schock}
\label{topic:schock-anaphylaktischer}
\index{Schock!Anaphylaktischer}
\index{Schock!allergischer}

\Definition{Schock aufgrund einer überschießenden
allergischen Reaktion -- "`Allergischer
Schock"'.}

%%%%%%%%%%%%%%%%%%%%%%%%%%%%%%%%%%%%%%%%%%%%%%%%%%%%%%%%%%%
% TODO : Section: Layout-Anpassung
\TODO{GABS}{2010-mm-yy}{Section: Layout-Anpassung notwendig.}{Dieser Abschnitt entspricht nicht dem Standard-Layout!
}
%%%%%%%%%%%%%%%%%%%%%%%%%%%%%%%%%%%%%%%%%%%%%%%%%%%%%%%%%%%

 Ausgelöst durch überschießende Reaktion der Körperabwehr (Immunreaktion) auf
bestimmte Stoffe (Antigene). 

% \begin{center}
%  [Warning: Image ignored] % Unhandled or unsupported graphics:
% %\includegraphics[width=2.641cm,height=3.837cm]{AASdev20090228-img70.jpg}
% \TempPicMissing{Allergischer Schock}
% \end{center}

Es werden sämtliche Gefäße maximal erweitert, zusätzlich
steigt die Durchlässigkeit der Gefäße
(Gefäßpermeabilität ${\uparrow}$) ${\rightarrow}$ 
Plasmaverlust in den Zwischengebwebsraum \Z{(Interstitium)}
${\rightarrow}$ RR fällt massiv ab!



\Layout{2}{Mögliche Auslöser}
{
% TODO VIT: Fließtext fehlt!
\begin{itemize}
    \item Pollen
    \item Medikamente, Antibiotika
    \item Kontrastmittel (für radiologische Untersuchungen)
    \item tierische Gifte (Bienenstich etc.)
    \item und vieles andere ...
\end{itemize}
}{}{}{}{}


\MarriedNG{\TxSymptome}
{
Tafel \ref{tab:anaphylaxie-stadien}.
}
{
Tafel \ref{tab:anaphylaxie-stadien}.
}
% TODO Flush erklären


\begin{table}[H]
\Captionof{table}{Stadien der
Anaphylaxie\label{tab:anaphylaxie-stadien}}
\begin{tabulary}{\linewidth}[H]{cL}
\TabHeading{\mbox{Stadium}}
 &
\\\toprule
\TabSub{I}
 &
Hautreaktion: evtl. Ausschlag, Rötung;
"`Quaddeln"', Flush, Ödeme

Allgemeinsymptome: Juckreiz, Zittern, Schwindel, Kopfschmerz

\\\addlinespace
\TabSub{II}

 &
RR ${\downarrow}$  , Tachykardie

Übelkeit, Erbrechen, Durchfall 

\\\addlinespace
\TabSub{III}
 &
Allgemeine Schocksymptome

Bronchospasmus
\\\addlinespace
\TabSub{IV}
 &
Atem-/Kreislaufstillstand

Schwerer Schock
\\\bottomrule
\end{tabulary}
\end{table}

\MassnahmenNG{Spezielle Maßnahmen}{Anaphylaktischer
Schock}{}{

\begin{itemize}
  \item Allgemein: \Es{Allgemeine Maßnahmen der
  Schockbekämpfung}
  (\prettyref{m:am-schock})
  \item Antigen  nach Möglichkeit entfernen!
  \item Lagerung: situationsgerecht.
  Grundsätzlich Beine hoch, bei
  Atemproblemen jedoch stark erhöhten
  Oberkörper erwägen.
\end{itemize}
}



\subsubsection{Septischer Schock und Toxischer Schock}
\index{Schock!Septischer}
\index{Schock!Toxischer}
\label{topic:schock-septischer}
\label{topic:schock-toxischer}
Siehe auch: Hygiene, \prettyref{topic:sepsis}.

\Definition{Schock infolge einer schweren, körperweiten Infektion
oder aufgrund eines, die Gefäße beeinflussenden, Giftstoffes
(Toxins).}

%%%%%%%%%%%%%%%%%%%%%%%%%%%%%%%%%%%%%%%%%%%%%%%%%%%%%%%%%%%
% TODO : Section: Layout-Anpassung
\TODO{GABS}{2010-mm-yy}{Section: Layout-Anpassung notwendig.}
{Dieser Abschnitt entspricht nicht dem Standard-Layout!}
%%%%%%%%%%%%%%%%%%%%%%%%%%%%%%%%%%%%%%%%%%%%%%%%%%%%%%%%%%%

Durch bakterielle Gefäßwandschädigung, temperaturbedingt (hohes
Fieber {\textgreater}\,38{\textcelsius}) sowie aufgrund von
Giftstoffen kommt es zu Flüssigkeitsaustritt in den Zwischen\-gewebs\-raum sowie
zu Gefäßweitstellung.

\Layout{2}{Ursachen}
{

%%%%%%%%%%%%%%%%%%%%%%%%%%%%%%%%%%%%%%%%%%%%%%%%%%%%%%%%%%%
% TODO : Fließtext/Slide fehlt!
\TODO{GABS}{2010-mm-yy}{Fließtext/Slide fehlt!}
   {Fließtext/Slide fehlt!}
%%%%%%%%%%%%%%%%%%%%%%%%%%%%%%%%%%%%%%%%%%%%%%%%%%%%%%%%%%%

Diverse Erreger und Toxine\index{Toxine!Toxischer Schock}

\begin{itemize}
  \item Langanhaltender Entzündungsherd der plötzlich Erreger
  "`ins Blut streut"' und "`explodiert"'
  \item Mit Keimen besiedelte Fremdkörper (diverse
  medizinische Sonden, Venflons, etc.)
  \item Tampons, die vergessen wurden.
  \item \ldots
\end{itemize}
}{}{}{}{}



\Layout{2}{\TxSymptome}
{

%%%%%%%%%%%%%%%%%%%%%%%%%%%%%%%%%%%%%%%%%%%%%%%%%%%%%%%%%%%
% TODO : Fließtext/Slide fehlt!
\TODO{GABS}{2010-mm-yy}{Fließtext/Slide fehlt!}{Fließtext/Slide fehlt!}
%%%%%%%%%%%%%%%%%%%%%%%%%%%%%%%%%%%%%%%%%%%%%%%%%%%%%%%%%%%

\begin{itemize}
  \item Schock
  \item hohes Fieber oder kein Fieber
  \item Haut heiß, trocken, rot oder bleich, grau,
  "`marmoriert"'
\end{itemize}
}{}{}{}{}








\MassnahmenNG{Spezielle Maßnahmen}{Septischer und Toxischer Schock}{}{
\begin{itemize}
  \item Allgemein: \Es{Allgemeine Maßnahmen der
  Schockbekämpfung} (\prettyref{m:am-schock})
  \item Lagerung: situationsgerecht, Beine hoch
  \item Schutz vor Unterkühlung/Überwärmung
  \item Aviso im Krankenhaus / Intensivstation
  ("`Überwachungsbett"')
  \item benötigt Antibiotika (im KH)
\end{itemize}
}

\subsubsection{Neurogener Schock}
\index{Schock!Neurogener}
\index{Schock!Spinaler}

\Definition{Schockzustand ausgelöst durch eine
Unterbrechung oder Störung der nervlichen Versorgung der
Kreislaufregulation. Es kommt zu inadäquater arterieller
und venöser Gefäßerweiterung.
\Z{(Vasodilatation)}}

\Layout{2}{Ursachen}
{
%%%%%%%%%%%%%%%%%%%%%%%%%%%%%%%%%%%%%%%%%%%%%%%%%%%%%%%%%%%
% TODO : Fließtext/Slide fehlt!
\TODO{GABS}{2010-mm-yy}{Fließtext/Slide fehlt!}{Fließtext/Slide fehlt!}
%%%%%%%%%%%%%%%%%%%%%%%%%%%%%%%%%%%%%%%%%%%%%%%%%%%%%%%%%%%

\begin{itemize}
    \item Neurologische Erkrankungen wie Querschnitt-Verletzung des
    Rückenmarks
    \item Die Gefäße bekommen keine Nerven-Impulse und kontrahieren
    sich nicht mehr.
    \item ${\rightarrow}$ Blut versackt
\end{itemize}
}{}{}{}{}





\MassnahmenNG{Spezielle Maßnahmen}{Neurogener Schock}{}{
\begin{itemize}
  \item Allgemein: \Es{Allgemeine Maßnahmen der
  Schockbekämpfung} (\prettyref{m:am-schock})
  \item Lagerung: Flachlagerung (WS-Verletzung ist Kontraindikation gg.
  Beine hoch!)
\end{itemize}
}




\vspace{2em}

%%%%%%%%%%%%%%%%%%%%%%%%%%%%%%%%%%%%%%%%%%%%%%%%%%%%%%%%%%%
%%%%%%%%%%%%%%%%%%%%%%%%%%%%%%%%%%%%%%%%%%%%%%%%%%%%%%%%%%%
%%%%%%%%%%%%%%%%%%%%%%%%%%%%%%%%%%%%%%%%%%%%%%%%%%%%%%%%%%%
%%%%%%%%%%%%%%%%%%%%%%%%%%%%%%%%%%%%%%%%%%%%%%%%%%%%%%%%%%%
\section{Temperaturregulation}
\label{topic:temperaturregulation}
\index{Temperatur!-regulation}

\A{KOCH: \#19: dient nur zur Info, wird nicht geprüft!}

Die normale Körper-Kerntemperatur beträgt
\Es{37\textcelsius}. Verschiedene Mechanismen sorgen für
Wärmeerzeugung und -abgabe:

%%%%%%%%%%%%%%%%%%%%%%%%%%%%%%%%%%%%%%%%%%%%%%%%%%%%%%%%%%%
% TODO : Section: Layout-Anpassung
\TODO{GABS}{2010-mm-yy}{Section: Layout-Anpassung notwendig.}{Dieser Abschnitt
entspricht nicht dem Standard-Layout! }
%%%%%%%%%%%%%%%%%%%%%%%%%%%%%%%%%%%%%%%%%%%%%%%%%%%%%%%%%%%

\begin{tabularx}{\linewidth}{XX}
\TabHeading{Wärmeproduktion durch}
&
\TabHeading{Wärmeabgabe über}
\\\addlinespace\midrule\addlinespace
\begin{itemize}
  \item Zellaktivität
  \item Stoffwechselvorgänge
  \item Bewegung (Muskelatmung)
\end{itemize}
&
\begin{itemize}
  \item Haut
  \item Schweiß
  
  \begin{itemize}
    \item Wirkung durch Verdunstungskälte
    \item ca. 0,5\,l\,/\,Tag
    \item bis max. 1,5\,\nicefrac{l}{h} bei
    schwerer körperlicher  Arbeit
  \end{itemize}
\end{itemize}
\end{tabularx}





%%%%%%%%%%%%%%%%%%%%%%%%%%%%%%%%%%%%%%%%%%%%%%%%%%%%%%%%%%%
%%%%%%%%%%%%%%%%%%%%%%%%%%%%%%%%%%%%%%%%%%%%%%%%%%%%%%%%%%%
%%%%%%%%%%%%%%%%%%%%%%%%%%%%%%%%%%%%%%%%%%%%%%%%%%%%%%%%%%%
%%%%%%%%%%%%%%%%%%%%%%%%%%%%%%%%%%%%%%%%%%%%%%%%%%%%%%%%%%%
\section{Stoffwechsel}
\label{topic:stoffwechsel}
\index{Stoffwechsel}
\index{Zuckerstoffwechsel}\index{Kohlenhydrate}\index{Fettstoffwechsel}
\index{Eiweißstoffwechsel}\index{Proteine}\index{Enzyme}

\Definition{Abläufe von Stoffumsetzung und 
Energiegewinnung, Nährstoffe werden in Wärme und Energie  umgewandelt}


\Layout{0}{\TxBeschreibung}
{
Der menschliche Körper hat 3 wesentliche Stoffwechselsysteme, die im zur
Verfügung stehen. Man unterscheidet den \Es{Zucker- bzw.
Kohlehydratstoffwechsel}, den \Es{Fettstoffwechsel} und den
\Es{Eiweißstoffwechsel}.

\subparagraph{Zuckerstoffwechsel} Zucker ist ein potenter \E{Energielieferant}
und kann vom Körper schnell umgesetzt werden. Er wird durch Aufspaltung von
komplexen \E{Kohlehydraten} (Stärke) gewonnen und zirkuliert in Form des
\Es{Blutzuckers} in der Blutbahn. Die Leber speichert eine Zuckerreserve von ca.
24\,h.

Das Hormon \T{Insulin}\label{topic:insulin}, welches in der Bauchspeicheldrüse
(Pankreas) produziert wird,  ermöglicht die \E{Aufnahme von Blutzucker in die
Zellen} und \E{senkt dabei den Blutzuckerspiegel}. 

\subparagraph{Fettstoffwechsel} Fette haben einen doppelt so hohen Energiegehalt
als Kohlenhydrate, können aber nicht so rasch vom Körper verwendet werden.
Sie dienen außerdem u.a. als Lösungsmittel für Vitamine, Bausteine der
Zellmembran und von Hormonen. \Z{(Weiters unterscheidet man
zwischen gesättigten und ungesättigten Fettsäuren.)}

\subparagraph{Eiweißstoffwechsel} Eiweiße (\T{Proteine}) sind vergleichbar mit
LEGO\texttrademark{}-Bausteinen: Sie sind aus Bausteinen zusammengebaut, den
\Z{Amminosäuren}. Zusammen gebaute Proteine können wieder "`\E{zerlegt}"'
werden, mit den entstehenden Amminosäuren können wiederum andere, \E{neue
Proteine gebaut} werden. 

Proteine sind in vielen Bereichen im menschlichen Körper unabkömmlich, sie
kommen zum Beispiel als \E{Enzyme} (\prettyref{topic:enzym}), \E{Gerinnungsfaktoren}
(\prettyref{topic:gerinnungsfaktoren}), Zellmembranbausteine u.a. zum Einsatz.
}{
\begin{itemize}
  \item Zuckerstoffwechsel (Kohlenhydrate)
  \begin{itemize}
    \item Energielieferant.
    \item Regulation durch Insulin.
  \end{itemize}
  \item Fettstoffwechsel \Z{(Fettsäuren, gesättigt und  ungesättigt)}
  \begin{itemize}
    \item Hoher Energiegehalt.
    \item Lösungsmittel für Vitamine, Bausteine der Zellmembran und von
    Hormonen.
  \end{itemize}
  \item Eiweißstoffwechsel (Proteine).
  \begin{itemize}
    \item Zellmembranbausteine.
    \item Enzyme, Gerinnungsfaktoren, \ldots
  \end{itemize}
\end{itemize}
}{}{}{}




\cite{AlbertsMolZellbiologie2,LoefflerBiochemie7}


\section{Immunsystem}
\label{topic:immunsystem}
% TODO Inhalt: neu: Immunsystem

\Layout{0}{\TxBeschreibung}
{
Das Immunsystem (Abwehrsystem) schützt den Körper vor Infektionskrankheiten.
Eine besonders bedeutende Rolle spielen die \Es{Leukozyten} (Weiße
Blutkörperchen, \prettyref{topic:leukozyten}). Sie ermöglichen die Abtötung und
den Abbau von Krankheitserregern im Körper. Fieber kann die Immunabwehr
unterstützen. 
}{
\begin{itemize}
  \item Schutz vor Infektionskrankheiten.
  \item Leukozyten.
  \item Fieber.
\end{itemize}
}{}{}{}


\Layout{0}{Störungen des Immunsystems}
{\subparagraph{Knochenmarkschädigende Therapien}Einige medizinische
Behandlungsmethoden verursachen als Nebenwirkung eine massive \Es{Schädigung des Knochenmarks}\index{Knochenmark!Schädigung} (z.\,B.
Bestrahlung, Chemotherapien bei Krebspatienten ("`onkologische Patienten"'),
etc.). Diese Patienten leiden neben einer Blutarmut auch an einem schwachen
Immunsystem, da nicht genügend Weiße Blutkörperchen gebildet werden. Sie können
daher sehr leicht schwerwiegende Infektionen bekommen.

Diese Patienten
sind im Krankentransport häufig anzutreffen.

\subparagraph{Allergie}
\label{topic:allergie}
\index{Allergie}
Das Immunsystem kann jedoch auch "`falsch" gegen
ungefährliche Stoffen aktiv werden, sozusagen als Irrlauf. Man spricht dann von
einer \T{Allergie}. Im leichtesten Fall kommt es bei Kontakt mit dem
auslösenden Stoff zu einer Schleimhautreizung ("`laufende Nase"'), Augenrötung,
tränen, einer lokalen Rötung bzw. einem Juckreiz. Je nach Stärke der Allergie
sowie Dauer und Intenensität des Kontakts kann es aber auch zu
lebensbedrohlichen Zuständen wie dem anaphylaktischen Schock kommen
(\prettyref{topic:schock-anaphylaktischer}). 
}{
\begin{itemize}
  \item Knochenmarkschädigende Therapien
  \begin{itemize}
    \item Schädigung des Knochenmarks durch med. Therapien (Chemotherapien,
    Bestrahlung, \ldots)
    \item Zu wenig eukozyten.
  \end{itemize}
  \item Allergien
  \begin{itemize}
    \item Fehlgeleitete Immunreaktion gegen ungefährliche Stoffe.
  \end{itemize}
\end{itemize}
}{}{}{}






\clearpage
\section{Flüssigkeit im Körper}
\index{Flüssigkeitshaushalt}
\index{Volumenhaushalt}

\Layout{2}{Verteilungsräume}
{
%%%%%%%%%%%%%%%%%%%%%%%%%%%%%%%%%%%%%%%%%%%%%%%%%%%%%%%%%%%
% TODO : Fließtext/Slide fehlt!
\TODO{GABS}{2010-mm-yy}{Fließtext/Slide fehlt!}{Fließtext/Slide fehlt!}
%%%%%%%%%%%%%%%%%%%%%%%%%%%%%%%%%%%%%%%%%%%%%%%%%%%%%%%%%%%

Das Körperwasser macht ca. 60\,\% des Körpergewichtes
aus und teilt sich auf:

\begin{itemize}
  \item \T{IntraZellularRaum} (IZR), 40\,\% KG
  \item \T{ExtraZellularRaum} (EZR), 20\,\% KG
  \begin{itemize}
    \item \index{Zwischengewebsraum}Zwischengewebsraum,
    $\sim$\,16\,\%: "`freier"'
    Raum zwischen den Zellen im Gewebe
    \item \index{Gefäßraum}Gefäßraum $\sim$\,4\,\%:
    Raum in den Blutgefäßen
  \end{itemize}
\end{itemize}
}{}{}{}{}


\Layout{0}{Trennung}
{
Die einzelnen Vertelungsräume haben Trennungen: Der Intrazellularraum
(\Abk{IZR}) wird vom Extrazellularraum (\Abk{EZR}) durch die Zellmembran
getrennt, der Zwischengewebsraum vom Gefäßraum durch die Blutgefäßwand.
}{
\begin{itemize}
    \item IZR -- EZR:    Zellmembran (semipermeable Membran, s.u.)
    \item Zwischengewebsraum -- Gefäßraum:  Blutgefäßwand
\end{itemize}
}{}{}{}


\Layout{20}{Atome, Ionen und Elektrolyte}
{
\index{Atome}
\index{Ionen}
\index{Elektrolyte}
\label{topic:atome}
\label{topic:ionen}
\label{topic:elektrolyte}
%%%%%%%%%%%%%%%%%%%%%%%%%%%%%%%%%%%%%%%%%%%%%%%%%%%%%%%%%%%
% TODO : Fließtext/Slide fehlt!
\TODO{GABS}{2010-mm-yy}{Fließtext/Slide fehlt!}{Fließtext/Slide fehlt!}
%%%%%%%%%%%%%%%%%%%%%%%%%%%%%%%%%%%%%%%%%%%%%%%%%%%%%%%%%%%


\T{Atome} sind die kleinsten Grundbestandteile von Marterie. \T{Ionen} sind
Atome die \E{elektrisch geladen} sind. Sie sind gut in Wasser löslich. Wenn sie
in Wasser gelöst sind heißen sie "`\T{Elektrolyte}"'.

\Fazit{Im Wasser gelöste, elektrisch geladene Teilchen
heißen \T{Elektrolyte}.}

Die wichtigsten Ionen sind Natrium (\ce{Na^+}), Kalium (\ce{K^+})
\Z{(positiv geladen)} und Chlorid (\ce{Cl^-}, \Z{negativ
geladen}).

\Z{Elektrolyte sind geladene Teilchen, die  bei Aufspaltung von Säuren, Laugen 
oder deren Salzen in wässriger Lösung  entstehen.  Man unterscheidet zwischen
positiv geladenen Teilchen, sog. \E{Kationen} (z.\,B. Kalium (\ce{K^+}), Natrium
(\ce{Na^+}), Magnesium (\ce{Mg^{++}}), Kalzium (\ce{Ca^{++}}), \ldots) und
negativ geladenen Teilchen, den \E{Anionen} (z.\,B. Chlorid (\ce{Cl^-})).

Diese Ionen sind unterscheidlich verteilt: Im \Abk{IZR} ist die
\ce{K^+}-Konzentration hoch, dagegen gibt es wenig \ce{Na^+}; Im 
\Abk{EZR} ist es umgekehrt: Es gibt viel \ce{Na^+} und wenig \ce{K^+}.
Spezielle Transportersysteme sorgend dafür. dass diese Ionen derartig verteilt
werden.}
}{

}
{./images/wikipedia-pd/Semipermeable_membrane.png}
{Gefäßraum und Zwischengewebsraum, getrennt durch die Blutgefäßwand. Es
schwimmt viel herum: Die grünen, gelben blauen und violetten Kugeln stellen
Ionen und Moleküle dar. Die roten Scheiben sind rote Blutkörperchen.}
{WmCo/pidalka44}


\Layout{0}{Funktion der Elektrolyte}{
Die Elektrolyte haben \E{zwei} wesentliche Funktionen: Einerseits haben sie
aufgrund der Osmose (s.u.) großen Einfluß auf die \Es{Flüssigkeitsverteilung}
im Körper, andererseits haben sie durch ihre elektrische Ladung
eine entsprechende \Es{elektrische Funktion}. Durch die elektrische
Ionenladung ist jede Körperzelle gegenüber ihrer Umgebung geladen. \Z{Würde
man winzig kleine Elektroden in eine Zelle und in ihre Umgebung stechen wäre
innen der Minuspol und außen der Pluspol.} Dies ist besonders für die
Erregungsleitung im Herzen, für die Nerven und den Muskel wichtig.
}{
\begin{itemize}
    \item \Es{Verschiebung von
    Wasser} zwischen den einzelnen Räumen \Z{(durch
    Osmose und Diffusion, siehe unten)}
    \item \Es{Elektrische Funktion} durch Ionenladung. Besonders
    wichtig für
    \begin{itemize}
        \item Herz (Erregungsleitung)
        \item Nerven
        \item Muskel
    \end{itemize}
\end{itemize}
}{}{}{}





%%%%%%%%%%%%%%%%%%%%%%%%%%%%%%%%%%%%%%%%%%%%%%%%%%%%%%%%%%%
%%%%%%%%%%%%%%%%%%%%%%%%%%%%%%%%%%%%%%%%%%%%%%%%%%%%%%%%%%%
%%%%%%%%%%%%%%%%%%%%%%%%%%%%%%%%%%%%%%%%%%%%%%%%%%%%%%%%%%%
\subsection{Wasser- und Elektrolythaushalt}
\index{Wasserhaushalt}
\index{Volumenhaushalt}
\index{Elektrolythaushalt}
\label{topic:wasser-elektrolythaushalt}

\Layout{0}{Exkurs}
{
Das Blutvolumen setzt sich zusammen aus dem \E{Plasma}, welches zum Großteil
aus Wasser besteht, und den \E{zellulären Bestandteile} (Blutkörperchen).

Der Anteil beträgt beim Erwachsenen 7-8\,\%, beim
Neugeborenes 8-9\,\% des Körpergewichts. \E{Verluste bis zu ca. 20\,\%} sind
prinzipiell ohne Funktionsbeeinträchtigung tolerierbar, allerdings
kommt es bald zu Schäden in der Endstrombahn. 

Osmose und Diffusion sorgen zusammen mit den Elektrolyten für die richtige
Verteilung der Flüssigkeit in den Verteilungsräumen und damit für die
Volumenregulation. Darüber hinaus steuern Hormone\ZFootnote{(Renin/Angiotensin/Aldosteron)} zusammen mit den Nieren die
planmäßige Ausscheidung von Flüssigkeit.
}{
\begin{itemize}
  \item Blutvolumen = Plasma(wasser) + zelluläre Bestandteile
  \item Erwachsener 7-8\,\%, Neugeborenes\,/\,Säugling 8-9\,\% des Körpergewichts
  \item Verluste bis 20\,\% verkraftbar.
  \item Osmose und Diffusion verteilen Wasser.
  \item Hormone \& Nieren sorgen für Ausscheidung.
\end{itemize}
}
{}{}{}









%%%%%%%%%%%%%%%%%%%%%%%%%%%%%%%%%%%%%%%%%%%%%%%%%%%%%%%%%%%
%%%%%%%%%%%%%%%%%%%%%%%%%%%%%%%%%%%%%%%%%%%%%%%%%%%%%%%%%%%
%%%%%%%%%%%%%%%%%%%%%%%%%%%%%%%%%%%%%%%%%%%%%%%%%%%%%%%%%%%
\subsection{Osmose und Diffusion}
\index{Osmose|Z}
\index{Diffusion|Z}
\label{topic:osmose}
\label{topic:diffusion}

\Definition{Verteilung von Wasser durch
Konzentrationsausgleich in verschiedene Räume, welche durch
halbdurchlässige Membranen voneinander getrennt sind.}

\AuthNotes{GAB: Osmose und Diffusion: Grundsätzlich wichtig und
relevant, allenfalls bessere Erklärung, aber bleibt in Normalschrift, nicht
heller. [GAB]

KOCH: Sollte alles grau werden}

\Layout{20}{Flüssigkeitsverschiebung durch Diffusion}
{
%%%%%%%%%%%%%%%%%%%%%%%%%%%%%%%%%%%%%%%%%%%%%%%%%%%%%%%%%%%
% TODO : Fließtext/Slide fehlt!
\TODO{GABS}{2010-mm-yy}{Fließtext/Slide fehlt!}{Fließtext/Slide fehlt!}
%%%%%%%%%%%%%%%%%%%%%%%%%%%%%%%%%%%%%%%%%%%%%%%%%%%%%%%%%%%

Geladene Teilchen können die Zellmembran nicht selbstständig
durchdringen, Wasser dagegen schon. Man nennt die Zellmembran daher
eine "`\index{semipermeable
Membran}halbdurchlässige\ZFootnote{semipermeable} 
Membran"'.

Die Natur möchte, dass in einer Flüssigkeit überall die
gleiche Konzentration von geladenen Teilchen herrscht. Da
die Zellmembran (und teilweise auch die Blutgefäßwand) nur das Wasser durchlässt,
nicht jedoch die geladenen Teilchen, verschiebt sich nur
das Wasser von einem Raum in den anderen. 

\Z{Man es sich so vorstellen, dass die
Elektrolyte einen "`Druck"' ausüben und das Wasser dazu bringen "`die Seiten zu
wechseln"'. Man nennt diesen Druck den
"`\index{osmotischer Druck}\T{osmotischen Druck}"'. 
Man kann auch sagen:
\Speech{"`Das Wasser folgt dem Konzentrationsgefälle."'}}
}{

}
% {./images/osmose.png}
{./images/wikipedia-pd/Osmotische_druk.png}
{Diffusion. }
{WmCo/Sander van der Molen, PD}





\Z{Daneben gibt es den "`onkotischen Druck"', welcher durch Eiweißbestandteile
ausgeübt wird.~\ZFootnote{Der durch Eiweißbestandteile hervorgerufene
onkotische Druck wird bei speziellen "`kolloidalen"' Infusionslösungen
ausgenützt. Sie enthalten große Eiweiße, welche die Blutgefäßwand
normalerweise nicht passieren können und folglich Wasser vom
Zwischengewebsraum in den Blutgefäßraum "`ziehen"'. Der verbreitetste Wirkstoff
ist die Hydroxyethylstärke (HAES), Beispiele für damit erzeugte
Infusionslösungen sind Elohaest{\texttrademark}, Voluven{\texttrademark} und
HyperHAES{\texttrademark}.}}

\Z{Wie kommen die Elektrolyte in die einzelnen Räume wenn doch nur Wasser
die Wand passieren kann? Dafür gibt es in der Zelle spezielle "`Pumpen"'
\Z{(Transportproteine)} die die Ionen hin und her
befördern. }

Die Steuerung der richtigen Verteilung der verschiedenen
Ionen ist recht kompliziert aber lebenswichtig. Eine "`falsche"' Verteilung der Elektrolyte kann schlimme Folgen haben (zB. Herzrhythmusstörungen). 

\Fazit{Die richtige Verteilung der Elektrolyte im Körper
ist lebenswichtig. Elektrolytstörungen können
lebensgefährlich sein. }

\Layout{0}{Häufige Ursachen von
Elektrolytstörungen}
{
\index{Elektrolytstörungen}
\label{topic:elektrolytstoerungen}
\label{topic:dialyse-elektrolyte}

\begin{itemize}
    \item Niereninsuffizienz (gestörte Ausscheidung)
    \begin{itemize}
        \item Besonders bei Dialysepatienten
    \end{itemize}
    \item Durchfall, Erbrechen (Verlust von Elektrolyten)
    \item Medikamente
    \item Vergiftung
    %\item \Z{Hormonentgleisungen}
\end{itemize}
}{
\begin{itemize}
    \item Niereninsuffizienz, Dialysepatienten
    \item Durchfall, Erbrechen (Verlust von Elektrolyten)
    \item Medikamente, Vergiftung
    %\item \Z{Hormonentgleisungen}
\end{itemize}
}{}{}{}




%%%%%%%%%%%%%%%%%%%%%%%%%%%%%%%%%%%%%%%%%%%%%%%%%%%%%%%%%%%
%%%%%%%%%%%%%%%%%%%%%%%%%%%%%%%%%%%%%%%%%%%%%%%%%%%%%%%%%%%
%%%%%%%%%%%%%%%%%%%%%%%%%%%%%%%%%%%%%%%%%%%%%%%%%%%%%%%%%%%
\subsection{Säure-/Basenhaushalt}
\index{Säure-Basen-Haushalt}
\label{topic:saeure-basen-haushalt}

\Layout{0}{Säuren und Laugen}
{
\index{Säure}
\index{Lauge}
\index{Base}
\T{Säuren} sind Stoffe, die 
Wasserstoffionen (\ce{H^+}) abgeben können, \T{Basen} sind 
Stoffe, die Wasserstoffionen (\ce{H^+}) abgeben können. Der Begriff
\T{Lauge} ist gleichbedeutend mit dem Begriff Base. 
}{
\begin{description}
  \item[\T{Säuren}] Stoffe, die
  Wasserstoffionen (\ce{H^+}) abgeben können
  \item[\T{Basen}\,/\,\T{Laugen}] Stoffe, die \ce{H^+} aufnehmen können
\end{description}
}{}{}{}



\Layout{2}{\T{pH}-Wert}
{
\index{pH-Wert}
%%%%%%%%%%%%%%%%%%%%%%%%%%%%%%%%%%%%%%%%%%%%%%%%%%%%%%%%%%%
% TODO : Fließtext/Slide fehlt!
\TODO{GABS}{2010-mm-yy}{Fließtext/Slide fehlt!}{Fließtext/Slide fehlt!}
%%%%%%%%%%%%%%%%%%%%%%%%%%%%%%%%%%%%%%%%%%%%%%%%%%%%%%%%%%%


\Definition{Der pH-Wert gibt an, wie \E{sauer} oder \E{basisch} eine
Lösung ist.}

pH\,{\textless}\,7 ist sauer, pH\,=\,7 ist neutral,
pH\,{\textgreater}\,7 ist basisch

\Z{Der pH-Wert beim Menschen ist normal ca. um 7,4.  Einen
pH unter 7,35  bezeichnet man als eine \E{Azidose}, einen
pH-Wert über 7,45  als \E{Alkalose}}

Der richtige pH-Wert ist lebenswichtig. Der Körper ist da sehr empfindlich, der
pH-Wert muss sich in engen Grenzen bewegen. Schon geringe
Abweichungen können große Folgen haben. Leider gibt es präklinisch normalerweise keine Möglichkeit
den pH-Wert zu messen, jedoch hat die Atmung bzw. Beatmung sowie Erkrankungen
 großen Einfluss auf ihn, siehe unten.

Bei einer Übersäuerung sprechen wir von einer
\index{Azidose}\T{Azidose}, wenn der Patient basisch wird, 
von einer \index{Alkalose}\T{Alkalose}.
% \Fazit{Bei einer \Es{Übersäuerung} sprechen wir von
% einer \T{Azidose}, wenn der Patient \Es{basisch} wird
% von einer \T{Alkalose}.}
}{}{}{}{}

\Layout{2}{Regulationsmechanismen}
{
%%%%%%%%%%%%%%%%%%%%%%%%%%%%%%%%%%%%%%%%%%%%%%%%%%%%%%%%%%%
% TODO : Fließtext/Slide fehlt!
\TODO{GABS}{2010-mm-yy}{Fließtext/Slide fehlt!}{Fließtext/Slide fehlt!}
%%%%%%%%%%%%%%%%%%%%%%%%%%%%%%%%%%%%%%%%%%%%%%%%%%%%%%%%%%%
Der Körper muß den pH-Wert in einem engen Bereich halten. Dazu gibt es
Regulationsmechanismen. Der wichtigste Regulationsmechanismus ist das
\Z{Bicarbonat-}Puffersystem im Blut.~\ZFootnote{Ein Puffer ist ein Stoff
welcher pH-Änderungen abfedert.}

Daneben gibt es auch noch andere Regulationsmechanismen, so reguliert auch die
Niere die Ausscheidung oder Wiederaufnhame von \ce{H^+}.

}{}{}{}{}

\Layout{0}{Puffersystem im Blut}
{
Der Mechanismus des \Z{Bicarbonat-}Puffers wandelt Säure \Z{(Bicarbonat und
\ce{H^+})} in Kohlendioxid \Z{\ce{CO2} und \ce{H2O}} um \E{und umgekehrt}.

\begin{equation}
\ce{$\text{\Z{\tiny Bikarbonat + }}$ $\text{Säure}$
        <->[\text{Gleichgewicht}] CO2 $\text{\Z{\tiny + Wasser}}$}
\end{equation}

\begin{equation}
\Z{\ce{HCO3^- + H^+ <->[\text{Gleichgewicht}] CO2 +
        H2O}}
\end{equation}        
        
Das \ce{CO2} wird normalerweise über die Lunge abgeatmet. Daraus folgt: 

\begin{enumerate}
  \item Atmet der Patient nicht (genug), staut
  sich das \ce{CO2} und \ce{H^+} (Säure) wird auf der anderen Seite nicht
  abgebaut.
  
  → \Es{Atmungsbedingte Übersäuerung}
  \item Atmet der Patient zu viel, wird zu viel
  \ce{CO2} abgeatmet und wird von \ce{H^+} und dem
  anderen Stoff nachgebildet 
  
  → \ce{H^+}
  fehlt → \Es{Atmungsbedingte Alkalose}
  \item Fällt zu viel \ce{H^+} (Säure) im
  Körper an (durch den Stoffwechsel, Stoffwechselbedingte Übersäuerung), wird
  es zu (Wasser und) \ce{CO2} umgewandelt, \Es{das \ce{CO2} wird
  abgeatmet.}
\end{enumerate}
}{
\begin{itemize}
  \item \ce{$\text{\Z{\tiny Bik. + }}$ $\text{Säure}$
  <->[\text{Gleichgewicht}] CO2 $\text{\Z{\tiny + Wasser}}$  }
  \item \ce{CO2} wird von Lunge abgeatmet.
  \begin{itemize}
    \item Zu wenig Atmung → \ce{CO2} staut
    sich → \ce{H^+} wird nicht abgebaut
    
    → \Es{Atmungsbedingte Übersäuerung}
    \item Zu viel Atmung → zu viel \ce{CO2} abgeatmet → \ce{CO2} wird
    nachgebildet → \ce{H^+}
    fehlt 
    
    → \Es{Atmungsbedingte Alkalose}
    \item Zu viel \ce{H^+} (Säure) im
    Körper → zu \ce{CO2} umgewandelt 
    
    → \Es{das \ce{CO2} wird
    abgeatmet.}
  \end{itemize}
\end{itemize}
}{}{}{}











\subsubsection{Atmungsbedingte Übersäuerung --
Respiratorische Azidose}
\label{bkm:RefHeading35233321}\label{topic:azidose-respiratorische}\index{Uebersäuerung@Übersäuerung}\index{Respiratorische
Azidose}durch Hypoventilation / verminderten  Gasaustausch:

\ce{CO2}  Abatmung vermindert ${\rightarrow}$ Anhäufung von
\ce{CO2} im Blut\footnote{\Z{Hyperkapnie}}, Ansäuerung des
Blutes


\Layout{2}{Ursachen}
{
%%%%%%%%%%%%%%%%%%%%%%%%%%%%%%%%%%%%%%%%%%%%%%%%%%%%%%%%%%%
% TODO : Fließtext/Slide fehlt!
\TODO{GABS}{2010-mm-yy}{Fließtext/Slide fehlt!}{Fließtext/Slide fehlt!}
%%%%%%%%%%%%%%%%%%%%%%%%%%%%%%%%%%%%%%%%%%%%%%%%%%%%%%%%%%%

\begin{itemize}
  \item Insuffiziente innere oder äußere Atmung
  \begin{itemize}
    \item Verlegung der Atemwege
    \item Atemfrequenz zu niedrig
    \item Atemzugsvolumen zu gering
    \item Sauerstoffangebot in der Einatemluft zu gering
    \item Gasaustausch beeinträchtigt
    \item Kreislauf zu schwach\footnote{Ist zwar nicht
    direkt atmungsbedingt, aber es kommt zu wenig
    Sauerstoff ins Gewebe.}
  \end{itemize}
\end{itemize}
}{}{}{}{}

\MassnahmenNG{Spezielle Maßnahmen}{Respiratorische Azidose}
{}{
\index{Alkalose!Respiratorische}

\begin{itemize}
    \item Behebung der Ursache
    \item Erhöhung des Atem-Minuten-Volumens
\end{itemize}
}






\subsubsection{Stoffwechselbedingte Übersäuerung
-- Metabolische Azidose} 
\label{topic:azidose-metabolische}
\index{Uebersäuerung@Übersäuerung}
\index{Azidose!Metabolische}
\index{Hyperventilation!bei metabolischer Azidose}

\Definition{Übersäuerung durch vermehrte Produktion von Säuren im Körper.}


\Layout{2}{Ursachen}
{
Zu der übermäßgen Säureproduktion kommt es vor allem bei einer
Kreislaufinsuffizienz und beim Schock.

Ein klassisches Krankheitsbild ist das Diabetische Koma. Auch hier werden saure
Stoffe gebildet \Z{Ketonkörper}, welche ein Ungleichgewicht verursachen. 
}{
\begin{itemize}
    \item Kreislaufinsuffizienz, Schock
    \item Diabetisches Koma
\end{itemize}
}{}{}{}


\Layout{2}{Gefahren}
{
Es kann zu einer Schädigung des Zellstoffwechsels und des Herzmuskels kommen.
}{
\begin{itemize}
    \item Schädigung des Zellstoffwechsels
    \item Schädigung des Herzmuskels
\end{itemize}
}{}{}{}

\MarriedNG{Symptome}
{
Es bestehen die Symptome der Grunderkrankung. Darüber hinaus fällt eine
\Es{tiefe, schnelle Atmung} zur Abatmung von \ce{CO2} (\E{Kußmaul'sche
Atmung}) auf.

}{
\begin{itemize}
    \item Symptome der Grunderkrankung
    \item Tiefe, schnelle Atmung (\E{Kußmaul'sche Atmung})
\end{itemize}
}


\subsubsection{Atmungsbedingte (respiratorische) Alkalose}
\label{topic:alkalose-respiratorische}
\index{Alkalose}
\index{Hyperventilation!respiratorische Alkalose}

%%%%%%%%%%%%%%%%%%%%%%%%%%%%%%%%%%%%%%%%%%%%%%%%%%%%%%%%%%%
% TODO : Section: Layout-Anpassung
\TODO{GABS}{2010-mm-yy}{Section: Layout-Anpassung notwendig.}
{Dieser Abschnitt entspricht nicht dem Standard-Layout!}
%%%%%%%%%%%%%%%%%%%%%%%%%%%%%%%%%%%%%%%%%%%%%%%%%%%%%%%%%%%

Durch die vermehrte Abatmung von \ce{CO2} wird der
pH-Wert des Blutes basisch und es kommt zu einem Elektrolytungleichgewicht.
Dies verursacht Symtome wie Bewusstseinsstörungen, Kribbeln in den Fingern,
Pfötchenstellung und tonischen Krämpfe. 

Eine ausführliche Beschreibung des \E{Hyperventilationssyndroms} findet sich
unter \prettyref{topic:hyperventilationstetanie}.

\subsubsection{Stoffwechselbedingte (metabolische) Alkalose}
\label{topic:alkalose-metabolische}
\index{Alkalose!Metabolische}
\index{Alkalose!Stoffwechselbedingte}

%%%%%%%%%%%%%%%%%%%%%%%%%%%%%%%%%%%%%%%%%%%%%%%%%%%%%%%%%%%
% TODO : Section: Layout-Anpassung
\TODO{GABS}{2010-mm-yy}{Section: Layout-Anpassung notwendig.}
   {Dieser Abschnitt entspricht nicht dem Standard-Layout!
}
%%%%%%%%%%%%%%%%%%%%%%%%%%%%%%%%%%%%%%%%%%%%%%%%%%%%%%%%%%%

Eine stoffwechselbedingte Alkalose  kommt v.A. durch den \E{Verlust von Säure},
z.\,B. durch anhaltendes Erbrechen, zu Stande.

\subsubsection*{Zusammenfassung}

\Fazit{Kurz gefasst:
\begin{itemize}
  \item Die \Es{richtige Verteilung} von Wasser im Körper ist lebenswichtig.
  \item Mittels Osmose und Diffusion verteilt der Körper das Wasser.
  \item Die \Es{Elektrolyte} haben einen großen Einfluß auf die
  Wasserverteilung: Nur wenn die Elektrolyte richtig verteilt sind, kann das
  Wasser richtig verteilt werden.
  \item Elektrolyte haben auch
  eine elektrische Funktion.
  \item Ein ausgeglichener Säure-Basen-Haushalt ist wichtig, damit chemische
  Abläufe im Körper richtig funktionieren können.
  \item Besonders die \Es{Atmung} und der \Es{Stoffwechsel} können dem
  Säure-Basen-Haushalt Streiche spielen:
  \begin{enumerate}
    \item Wer zu wenig atmet wird sauer.
    \item Wer zu viel atmet wird basisch und spürt ein Kribbeln in den Fingern.
    \item Wer Probleme mit dem Stoffwechsel hat kann auch
    sauer oder basisch werden.
  \end{enumerate}
\end{itemize}

}

\cite{LoefflerBiochemie7}




\clearpage
%%%%%%%%%%%%%%%%%%%%%%%%%%%%%%%%%%%%%%%%%%%%%%%%%%%%%%%%%%%
%%%%%%%%%%%%%%%%%%%%%%%%%%%%%%%%%%%%%%%%%%%%%%%%%%%%%%%%%%%
%%%%%%%%%%%%%%%%%%%%%%%%%%%%%%%%%%%%%%%%%%%%%%%%%%%%%%%%%%%
%%%%%%%%%%%%%%%%%%%%%%%%%%%%%%%%%%%%%%%%%%%%%%%%%%%%%%%%%%%
\section{Tod}
\index{Tod}
\label{topic:tod}

Der Tod ist eine der grundlegensten und unausweichlichsten
Vitalfunktionen überhaupt.  



\paragraph{Klinischer Tod vs. biologischer Tod}
\index{Tod!klinischer}\index{Tod!biologischer}

\begin{center}
\begin{tabularx}{\linewidth}{XX}
\toprule
\TabHeading{Klinischer Tod}
 &
\TabHeading{Biologischer Tod}
\\\midrule
keine Atmung

kein Puls

Reanimation muss innerhalb der ersten  Minuten erfolgen

(dann unter Umständen) reversibel
 &
Organtod des Gehirns (irreversibel)

gleichzusetzen mit Tod des Individuums
\\\bottomrule
\end{tabularx}
\end{center}

% \paragraph{Wann ist es wirklich zu spät?}



\AuthNotes{Todeszeichen: Nicht wie von Roman gefordert in
Tabelle mit mit sicheren und unsicheren Todeszeichen umgebaut, halte ich für unübersichtlich
und befürchte Verwechslungsgefahr mit
Knochenbruchzeichen.[GAB]}





\Layout{2}{Zeichen des Todes und Leichenzeichen}
{
Grundsätzlich gibt es 3 Zeichen die als \Es{sichere
Todeszeichen} gelten\cite{ForMed2}:

\begin{enumerate}
    \item \Es{Bildung von blauvioletten Totenflecken} durch
    Absackung des Blutes an die untenliegenden Stellen.
    \item \Es{Nicht mit dem Leben vereinbare Verletzungen}
    wie zum Beispiel eine Abtrennung des Kopfes vom Körper.
    \item \Es{Fäulniserscheinungen}
\end{enumerate}

Die ersten beiden Punkte sind \E{frühe Leichenerscheinungen}, da sie relativ
rasch auftreten. Andere frühe Leichenerscheinungen sind die \E{Abkühlung der
Leiche} und die \E{Totenstarre}. Längere Zeit nach dem Tod treten die \E{späten
Leichenerscheinungen} auf, dazu gehören \E{Fäulnis und Verwesung},
\E{Mumifizierung}, \E{Tierfraß} und \E{Skelettierung}.

Die Feststellung der Todeszeichen und Leichenerscheinungen
ist jedoch oft nicht so einfach, die Anzeichen können oft
auch mit anderen Sachen verwechselt werden, z.\,B.
Totenflecke mit Blutergüssen oder die Totenstarre mit einem
tonischen Krampf.
}{}{}{}{}




\Fazit{Die Todesfeststellung erfolgt grundsätzlich nur durch einen
Arzt!}\index{Todesfeststellung}

\Cave{Nobody is dead until warm and dead!

Niemand wird für tot erklärt, bevor er nicht erwärmt
wurde.}

%\vspace{1em}
%%%%%%%%%%%%%%%%%%%%%%%%%%%%%%%%%%%%%%%%%%%%%%%%%%%%%%%%%%%
%%%%%%%%%%%%%%%%%%%%%%%%%%%%%%%%%%%%%%%%%%%%%%%%%%%%%%%%%%%
%%%%%%%%%%%%%%%%%%%%%%%%%%%%%%%%%%%%%%%%%%%%%%%%%%%%%%%%%%%
\subsection{Unterlassung der Reanimation}
\label{topic:reanimation-unterlassung}


              

\begin{gWideParEnv}
\fcolorbox{ubuntured}{White}{
\begin{minipage}[t]{\linewidth - {2\fboxrule} - {2\fboxsep}}
\Ca{Unter den folgenden Voraussetzungen darf (oder muss) eine
Reanimation unterlassen
werden:\index{Reanimation!Unterlassung
der}\index{Wiederbelebung|see{Reanimation}}}
\begin{multicols}{2}              
\begin{itemize}
  \item \Es{Verwesung}
  \item \Es{Fäulnis}
  \item \Es{Mumifizierung}
  \item \Es{Skelettierung}
  \item \Es{Verletzungen, die nicht mit dem Leben
  vereinbar sind}
  \item \Es{Tierfraß}\footnote{Unter
  \E{"`Tierfraß"'} versteht man den Verzehr eines Kadavers
  durch Kadaverfresser. Ein frischer Hundebiss oder
  Ähnliches fällt \underbar{nicht} darunter. Maden können
  sich im Gewebe von Lebenden einnisten und sind daher
  auch kein sicheres Todeszeichen.}
  \item Verbindliche % oder Beachtliche
  % TODO Beachtliche Patientenverfügung: keine Rea ???
  % BEachtliche erstmals ausgeklammert, STEU und GABS befürchten Verwirrung
  % der Leser. Evtl. eine gute Ausformulierung nötig, evtl. als Fußnote?
  \Es{Patientenverfügung}\index{Patientenverfügung!Unterlassung der
  Reanimation} (\prettyref{topic:patientenverfuegung})
  \item Abbruch der Reanimationsmaßnahmen aufgrund
  von \Es{Erschöpfung}
  \item Unvereinbarkeit mit
  \Es{Selbstschutz}\footnote{Die Unvereinbarkeit mit dem Selbstschutz ist kritisch zu
  betrachten. Es müssen auf jeden Fall alle Maßnahmen
  getroffen werden die ohne Gefährdung des Selbstschutzes
  möglich sind. So wird zum Beispiel für Ersthelfer
  empfohlen nur eine Herzdruckmassage durchzuführen wenn
  eine Mund-zu-Mund-Beatmung nicht zumutbar ist.}
\end{itemize}

\end{multicols}
\end{minipage}
}
\end{gWideParEnv}








\AuthNotes{GABG: gesicherte bösartige Erkrankung?

GABS: Was gilt als gesichert? Was gilt als bösartig? Eine
bösartige Grunderkrankung gilt nicht per se als Argument um
nicht zu reanimieren $\rightarrow$ Patientenverfügung,
mutmaßlicher Patientenwille. Prostatakarzinom und Infarkt?}


%O2:
%\nocite{Morris2001,Milross1989,Boumphrey2003}
%\nocite{BoehmerTaschRett6,Fauci2008,LPN3,LippertAnatomie7,StriebelAn7,DRAn3}


%\putbib

\nocite{amls3,emergency9,phtls6,RedelsteinerHRNS1}

\ChapterBibliography




